\documentclass[10pt]{article}

\usepackage[margin=1in]{geometry}
\usepackage{amsmath,amssymb,amsthm,mathtools,bm}
\usepackage{booktabs,array}
\usepackage{microtype}
\usepackage{xcolor}
\usepackage{enumitem}
\usepackage{hyperref}
\usepackage{url}

\hypersetup{
  colorlinks=true,
  linkcolor=blue!55!black,
  citecolor=green!45!black,
  urlcolor=blue!65!black
}

\allowdisplaybreaks
\setlist[itemize]{leftmargin=*,topsep=3pt,itemsep=2pt}
\setlist[enumerate]{leftmargin=*,topsep=3pt,itemsep=2pt}

\newtheorem{theorem}{Theorem}[section]
\newtheorem{proposition}[theorem]{Proposition}
\newtheorem{lemma}[theorem]{Lemma}
\newtheorem{corollary}[theorem]{Corollary}
\newtheorem{assumption}[theorem]{Assumption}
\theoremstyle{definition}
\newtheorem{definition}[theorem]{Definition}
\newtheorem{example}[theorem]{Example}
\theoremstyle{remark}
\newtheorem{remark}[theorem]{Remark}

\newcommand{\R}{\mathbb{R}}
\newcommand{\N}{\mathbb{N}}
\newcommand{\E}{\mathbb{E}}
\newcommand{\Pp}{\mathbb{P}}
\newcommand{\KL}{\mathrm{KL}}
\newcommand{\TV}{\mathrm{TV}}
\newcommand{\dd}{\,\mathrm{d}}
\newcommand{\ind}{\mathbf{1}}
\newcommand{\emptyconfig}{\varnothing_M}
\newcommand{\state}{\mathsf{X}}
\newcommand{\pdata}{P_{\mathrm{data}}}
\newcommand{\splitmap}{\mathcal{S}}
\newcommand{\mergemap}{\mathcal{C}}
\newcommand{\deathmap}{\mathcal{D}}
\newcommand{\birthmap}{\mathcal{B}}

\newenvironment{algorithmblock}[1]
{\par\medskip\noindent\begin{minipage}{\linewidth}\hrule\medskip
 \noindent\textbf{#1}\par\smallskip\small}
{\par\medskip\hrule\end{minipage}\par\medskip}

\title{\textbf{Mass-Conserving Coagulation--Fragmentation Generative Models}\\
Finite-Horizon Time Reversal on Variable-Size Weighted Sets}
\author{Anonymous authors}
\date{Draft: theory and algorithms complete; experiments not yet run}

\begin{document}

\maketitle

\begin{abstract}
Many data objects are unordered collections with a variable number of components and a meaningful positive additive quantity, such as mass, area, energy, or probability weight.  Standard diffusion models usually keep the number of components fixed, while a prominent trans-dimensional alternative removes and inserts one component at a time without an additive pair-merge constraint.  We introduce a different construction based on conservative binary fragmentation and coagulation.  The data-to-reference process first fragments components and moves component mass into an explicit reservoir.  It then removes every remaining component, so that the terminal state is exactly an all-reservoir state.  The finite-horizon reverse process starts from this simple state and performs nucleation, growth, and pairwise coagulation.

The paper makes the construction precise on an unordered, variable-dimensional state space.  We define compatible reference measures, conservative split and merge maps, and a nonnegative reservoir flow.  We derive the reverse jump kernel from probability flux, including the mass Jacobian and the cancellation induced by a consistent unordered-pair convention.  We then specialize a path-space evidence lower bound to this conservative marked-event process.  Its trainable part is the negative log likelihood of reversed forward trajectories under a leaf-count, marked-event model.  This gives a simple simulation-based learning algorithm that does not require evaluating the unknown marginal density ratio appearing in the exact reverse rates.  We prove terminal absorption, non-explosion, exact mass conservation, permutation invariance, Fisher consistency under stated model-class assumptions, and an endpoint error bound.  We also give complete forward-simulation, training, and reverse-sampling algorithms.  An optional extension replaces uniform evaporation with a spatial control field motivated by clearing a fogged mirror one region at a time.  The experimental protocol is specified, but no experimental results are reported in this draft.
\end{abstract}

\section{Introduction}

A set-valued observation can be written informally as
\begin{equation}
  X=\left\{(m_i,z_i)\right\}_{i=1}^{N},
  \qquad m_i>0,\quad z_i\in\R^q,
  \label{eq:intro-set}
\end{equation}
where the braces are understood as an unordered multiset, the order of the
elements has no meaning, and the cardinality $N$ may change from one example
to another.  The scalar $m_i$ is an additive quantity.  Depending on the application, it may represent physical mass, energy, area, intensity, or a normalized weight.  The vector $z_i$ stores position and other continuous features.

We study generative models for such observations under the exact constraint
\begin{equation}
  v+\sum_{i=1}^{N}m_i=M.
  \label{eq:mass-constraint-intro}
\end{equation}
Here $M>0$ is fixed or supplied as conditioning information, and $v\geq 0$ is an uncondensed reservoir.  Introducing $v$ is important.  It lets a component grow or evaporate without violating the additive constraint.

The physical picture is condensation, but the method is not obtained by replacing a Gaussian noise schedule with the Hertz--Knudsen equation.  The useful structure is instead the change in topology:
\[
\text{reservoir}
\longrightarrow \text{nucleation}
\longrightarrow \text{growth}
\longrightarrow \text{pairwise coagulation}
\longrightarrow \text{sample}.
\]
Training uses the opposite direction.  A prescribed forward process fragments data components, transfers their mass to the reservoir, and finally deletes them.  The learned model follows the time-reversed trajectories.

The motivating mirror observation also contains a second idea.  A hair dryer
does not clear the whole mirror uniformly: it acts on one small region, the
user observes what remains covered, and then moves the dryer.  Section~\ref{sec:local-control}
turns this observation into an optional spatially controlled evaporation
field.  The base model remains the uniform-control special case.

This paper has five main contributions.
\begin{enumerate}
  \item We define a conservative fragmentation--evaporation process on finite unordered weighted configurations.  A two-stage schedule makes the all-reservoir terminal state exact rather than approximate.
  \item We specialize finite-horizon Markov-process reversal to this state
  space.  The resulting reverse events are nucleation and pairwise
  coagulation.  The coordinate derivation handles both the mass-dependent
  Jacobian and unordered-pair counting.
  \item We specialize a path-space evidence lower bound to the conservative
  marked-event process and obtain a tractable point-process loss.  Forward
  paths provide supervised reverse events, while a survival integral trains
  merge timing and stopping.
  \item We give implementable neural parameterizations and event-driven algorithms with exact conservation, positivity, permutation invariance, and explicit non-explosion conditions.
  \item We formulate a locally controlled evaporation extension, prove that
  it remains conservative and positive, and explain its conditional reverse,
  flow Jacobian, feedback policy, and implementation.
\end{enumerate}

The contribution is deliberately narrower than ``a condensation diffusion model.''  Variable-dimensional generation and general jump-process time reversal already exist.  Our focus is the complete conservative one-to-two/two-to-one construction on unordered weighted sets.

\section{Related work}

\paragraph{Corruption and reversal.}
Diffusion models learn to reverse a prescribed corruption process
\cite{sohldickstein2015deep,ho2020denoising,song2021score}.  Their connection to nonequilibrium thermodynamics is therefore foundational rather than new.  Cold Diffusion showed that useful models can invert degradations other than Gaussian noising \cite{bansal2023cold}.  Consequently, calling a degradation ``evaporation'' is not by itself a methodological contribution.

\paragraph{Trans-dimensional generation.}
The closest generative framework is the Jump Diffusion Model of
Campbell et al.\ \cite{campbell2023transdimensional}.  It deletes individual components in the forward direction and learns state-conditioned single-component insertion in reverse, together with a trans-dimensional reversal formula and an evidence lower bound.  Learning to Jump uses thinning and thickening for latent counts \cite{chen2023learning}.  The 2026 preprint Existence-Field Diffusion instead represents variable cardinality by continuous existence variables \cite{pan2026existence}.  We do not claim that variable-dimensional generation is new.  Our distinction is a conservative binary split whose reverse must select and merge an interacting pair.

\paragraph{Structured expansion and branching.}
Iterative Local Expansion reverses graph coarsening by repeatedly splitting coarse nodes and refining local structure \cite{bergmeister2024efficient}.  The 2026 preprint Hierarchical Discrete Flow Matching also uses coarse-to-fine graph expansion \cite{boget2026hierarchical}.  In particle physics, JUNIPR generates particle sets through conservation-aware binary branchings \cite{andreassen2019junipr}.  These papers show that structured splitting is not itself new.  In our construction, the genealogy is sampled by the prescribed corruption rather than fixed by a preprocessing-induced clustering tree.

Dimension-changing maps and their Jacobians also have a long history in
reversible-jump Markov chain Monte Carlo \cite{green1995reversible}.  Unlike
that stationary MCMC setting, we study a learned finite-horizon reversal.

\paragraph{Coagulation, fragmentation, and reversal.}
Coagulation--fragmentation processes have a large mathematical literature
\cite{aldous1999deterministic,norris1999smoluchowski,bertoin2006random}.  General time-reversal results also exist for diffusions and Markov processes with jumps
\cite{anderson1982reverse,conforti2022time}.  We therefore do not present the probability-flux identity as a new general theorem.  The technical work is its specialization to conservative unordered configurations, including compatible measures, event coordinates, counting factors, a trainable objective, and a practical sampler.

\section{State space and notation}
\label{sec:state}

\subsection{Weighted configurations}

Fix a total mass $M>0$, a feature dimension $q\geq 1$, and a finite
cardinality cap $N_{\max}\geq1$.  The feature space is
$\mathcal Z=\R^q$, equipped with its Borel sets and $q$-dimensional
Lebesgue measure $\lambda_q$.  We use the shorthand
\begin{equation}
  \lambda_q^{\otimes n}(\dd(z_1,\ldots,z_n))
  \equiv\dd z_{1:n}
  :=\prod_{i=1}^n\lambda_q(\dd z_i).
  \label{eq:feature-reference-measure}
\end{equation}
Thus every density in this paper that contains $\dd z_{1:n}$ is a density
for continuous features relative to product Lebesgue measure.  Discrete or
mixed features require a different product reference measure and are not
covered by the displayed coordinate formulas.
All state spaces and reference measures below are conditional on the chosen
$M$; this dependence is suppressed in the notation $\state_n$ and $\mu_n$.

A state with $n$ components is
\begin{equation}
  X=\left(v;\{(m_i,z_i)\}_{i=1}^{n}\right),
  \qquad
  v\geq 0,\quad m_i>0,\quad z_i\in\R^q,
  \qquad v+\sum_{i=1}^{n}m_i=M.
  \label{eq:state}
\end{equation}
The braces in \eqref{eq:state} denote an unordered multiset, not an ordinary
set: identical pairs may occur with multiplicity.  Formally, a state is an
equivalence class of labeled tuples under permutations of the $n$ indices.
Two states are equal when one can be obtained from the other by such a
permutation.  We write $N(X)=n$ and
\begin{equation}
  \mathcal{M}(X)=v+\sum_{i=1}^{N(X)}m_i.
\end{equation}
The all-reservoir state is
\begin{equation}
  \emptyconfig=(M;\varnothing).
\end{equation}

It is sometimes useful to view the components as the finite measure
\begin{equation}
  \eta_X=\sum_{i=1}^{N(X)}m_i\delta_{z_i}.
  \label{eq:finite-measure}
\end{equation}
Then $\eta_X(\mathcal Z)=M-v$, and relabeling the components does not change
$\eta_X$.  The symbol $\eta_X$ describes one state; it is not the reference
measure used to integrate over all states.  If two components have exactly
the same feature $z$, $\eta_X$ combines their masses and no longer records
their multiplicity.  The unordered multiset in \eqref{eq:state} is therefore
the primary representation.  Under the continuous feature densities used
below, exact feature collisions have measure zero.

Let $\state_n$ be the space of states with $n$ components and let
\begin{equation}
  \state=\bigsqcup_{n=0}^{N_{\max}}\state_n
  \label{eq:disjoint-union}
\end{equation}
be their disjoint union.  For the finite-horizon process introduced in
Section~\ref{sec:forward}, $t\in[0,T]$ denotes forward time and
$X_t\in\state$ is the random state at time $t$.  Set
$N_t=N(X_t)$.  On the event $\{N_t=n\}$, we have $X_t\in\state_n$.
The process $N_t$ is piecewise constant and changes only at a split or
deletion event in forward time, or at a birth or merge event in reverse
time.  The cap is part of the model: a forward split is
disabled when $N(X)=N_{\max}$.  It prevents an artificial explosion of
ever smaller fragments and makes the reverse leaf budget finite.  A large
$N_{\max}$ has no practical effect when data and simulated paths remain far
below it.

\paragraph{Data law, corruption law, and model law.}
These three laws have different roles.  The raw data are sampled from an
unknown law on unordered weighted sets; we do not assume that nature
generated them by fragmentation.  The forward fragmentation--evaporation
process is a user-chosen conditional corruption law used to construct
training histories.  The learned reverse process is the generative model:
it samples a leaf budget, birth times and marks, deterministic growth, and
merge events, thereby inducing a joint endpoint law for masses and
features.  In particular, there is no separate uniform law for $z_{1:n}$,
and no independence among components, masses, and features is assumed.

\paragraph{Reservoir dequantization.}
Suppose a nonempty raw observation has normalized positive weights
\begin{equation}
  w=(w_1,\ldots,w_n)\in\Delta_{n-1}^{\circ}
  :=\left\{w\in(0,1)^n:\sum_{i=1}^n w_i=1\right\}.
  \label{eq:open-simplex}
\end{equation}
Only $w_{1:n-1}$ are free, because
$w_n=1-\sum_{i<n}w_i$.  We assume the raw weights have a density in these
$n-1$ simplex coordinates and that the features have a density relative to
\eqref{eq:feature-reference-measure}.  The corresponding unordered raw
reference measure is
\begin{equation}
  \nu_n^{\mathrm{raw}}(\dd w\,\dd z)
  =\frac{1}{n!}\ind\{w\in\Delta_{n-1}^{\circ}\}
   \dd w_{1:n-1}\,\dd z_{1:n}.
  \label{eq:raw-reference-measure}
\end{equation}
The displayed coordinates use any labeled representative, with $w_n$
eliminated.  Permuting the representative only changes to another simplex
chart with absolute Jacobian one, and the factor $1/n!$ makes the resulting
unordered density label-independent.  The case $n=1$ has no free $w$
coordinate; its simplex is a single point.  The empty outcome has no
normalized weight vector and is handled separately in $\state_0$.

A fixed reservoir fraction $\xi_0$ would give
$\sum_i m_i=M(1-\xi_0)$.  The mass vector would therefore lie on an
$(n-1)$-dimensional hyperplane inside the $n$-dimensional mass domain.  That
hyperplane has zero $\dd m_{1:n}$ measure, so the augmented data law would
be singular relative to the state reference measure defined below.  A
continuous random reservoir fraction supplies the missing radial dimension.
Specifically, draw
\begin{equation}
  \xi\sim q_\epsilon,
  \qquad
  v=\xi M,
  \qquad
  m_i=(1-\xi)Mw_i,
  \label{eq:reservoir-dequantization}
\end{equation}
where $q_\epsilon$ is a positive smooth density, relative to $\dd\xi$, on
$I_\epsilon=[\epsilon_{\min},\epsilon_{\max}]\subset(0,1)$.  In this paper
$q_\epsilon$ is fixed and independent of $(w,z)$; a data-dependent
dequantization density could be used if its conditional log density were
included in the bound.  The resulting state has $v>0$ and
$\sum_i m_i<M$, and its law can have a full-dimensional density.  The
dequantization density and its Jacobian must be included when a numerical
data-likelihood bound is reported.  In
coordinates $(w_1,\ldots,w_{n-1},\xi)\mapsto(m_1,\ldots,m_n)$, with
$w_n=1-\sum_{i<n}w_i$, the absolute Jacobian determinant (see
Appendix~\ref{app:dequantization-derivation}) is
\begin{equation}
  M^n(1-\xi)^{n-1}.
  \label{eq:dequantization-jacobian}
\end{equation}
At the output, no auxiliary $\xi$ needs to be stored.  For a nonempty state,
let $c=\sum_jm_j=M-v>0$ and decode
\begin{equation}
  \widehat w_i=\frac{m_i}{c}=\frac{m_i}{\sum_jm_j}.
  \label{eq:normalized-decoder}
\end{equation}
If we define the state-implied fraction $\widehat\xi=v/M$, then
$c=M(1-\widehat\xi)$, so this is also
$\widehat w_i=m_i/[M(1-\widehat\xi)]$.  The empty state is a separate
discrete outcome.  If an application observes a genuine reservoir, no
augmentation or normalization decoder is needed provided that its joint law
is supported on $v>0$ and is dominated by the interior reference measure
below.  An atom at $v=0$ requires an additional boundary reference measure
or an analogous dequantization.

\begin{proposition}[The normalized encoder and decoder agree]
\label{prop:encode-decode}
For every nonempty state, the numbers in
\eqref{eq:normalized-decoder} are positive and sum to one.  If a raw
normalized set is encoded by \eqref{eq:reservoir-dequantization}, decoding
it recovers its original weights exactly.  Conversely, decoding any
nonempty state and re-encoding it with $\widehat\xi=v/M$ recovers all of its
masses and its reservoir; the features are copied unchanged in both maps.
\end{proposition}

\begin{proof}
Because every $m_i>0$ and $c=\sum_jm_j>0$, each $\widehat w_i>0$, and
$\sum_i\widehat w_i=c/c=1$.  For an encoded datum,
$c=\sum_iM(1-\xi)w_i=M(1-\xi)$, hence
$m_i/c=w_i$.  Conversely, with $\widehat\xi=v/M$,
\[
  M(1-\widehat\xi)\widehat w_i
  =(M-v)\frac{m_i}{M-v}=m_i,
  \qquad M\widehat\xi=v.
\]
Thus the two maps are inverses on every nonempty stratum.
\end{proof}

For completeness, the raw normalized-data likelihood obeys the standard
dequantization bound
\begin{align}
  \log p_\theta^{\mathrm{raw}}(w_{1:n-1},z_{1:n})
  \geq \E_{\xi\sim q_\epsilon}\big[&
    \log p_{\theta,n}(X(\xi))
    +n\log M+(n-1)\log(1-\xi)
    -\log q_\epsilon(\xi)\big].
  \label{eq:dequantization-bound}
\end{align}
Here $p_\theta^{\mathrm{raw}}$ is a density relative to
$\nu_n^{\mathrm{raw}}$, whereas $p_\theta$ is a state density relative to
the measure below.  More explicitly, if $P_\theta$ is the endpoint law of
the learned reverse process, then $p_{\theta,n}$ is its density on
$\state_n$.  The bound follows by marginalizing the state-implied
reservoir fraction over $(0,1)$, restricting the integral to
$I_\epsilon$, inserting $q_\epsilon/q_\epsilon$, and applying Jensen's
inequality.  The path ELBO
below can be substituted for $\log p_{\theta,n}(X(\xi))$ to obtain one combined
lower bound.
Appendix~\ref{app:dequantization-derivation} gives the complete derivations
of \eqref{eq:dequantization-jacobian} and
\eqref{eq:dequantization-bound}.
Throughout the theory below, $\pdata$ denotes the law of this augmented
state $X$ (or the directly observed state when a reservoir is available).

\subsection{Reference measure}

An unordered density requires a counting convention.  Let
$\mathfrak q_n$ map a labeled coordinate tuple $(m_{1:n},z_{1:n})$ to its unordered
state, with $v=M-\sum_i m_i$.  We define $\mu_n$ by requiring that, for
every nonnegative measurable function $f$, where
$\dd m_{1:n}:=\prod_{i=1}^n\dd m_i$ is $n$-dimensional Lebesgue volume,
\begin{align}
  \int_{\state_n}f(X)\mu_n(\dd X)
  =\frac{1}{n!}\int
  &f\!\left(\mathfrak q_n(m_{1:n},z_{1:n})\right)
  \ind\!\left\{m_i>0,\ \sum_{i=1}^{n}m_i<M\right\}
  \dd m_{1:n}\,\dd z_{1:n}.
  \label{eq:reference-measure}
\end{align}
We use the corresponding shorthand
$\mu_n(\dd X)=n!^{-1}\ind\{m_i>0,\sum_i m_i<M\}
\dd m_{1:n}\dd z_{1:n}$.

For fixed $M$, the mass-coordinate domain is the open scaled simplex
\begin{equation}
  D_n(M)=\left\{m\in(0,\infty)^n:\sum_{i=1}^n m_i<M\right\}.
  \label{eq:mass-domain}
\end{equation}
Each individual mass lies in $(0,M)$, but the joint domain is not the full
box $(0,M)^n$ because the masses share the same budget.  Lebesgue volume on
$D_n(M)$ is a reference measure, not a uniform probability distribution.

The condition $m_i>0$ repeats the definition of $\state_n$ in ambient
Euclidean coordinates so that the domain of integration is explicit.  A
zero-mass element is absent and belongs to a lower-cardinality stratum.  The
strict inequality selects the interior $v>0$.  We keep the boundary $v=0$
in the state space for event-map and flow statements, but it has zero
$\mu_n$ measure; a law with an atom on that boundary would need an additional
singular measure component.

The factor $1/n!$ is needed because a generic unordered configuration has
$n!$ labeled coordinate representations.  Integrating a symmetric function
over all labeled coordinates without this factor would count the same state
$n!$ times.  The state-specific finite measure $\eta_X$ in
\eqref{eq:finite-measure} does not solve this counting problem: it represents
one configuration, whereas $\mu_n$ defines volume across all configurations.
For example, when $n=2$, the labeled tuples
$((m_1,z_1),(m_2,z_2))$ and
$((m_2,z_2),(m_1,z_1))$ describe the same unordered state.  The factor
$1/2!$ counts that state once.  Equivalently, one could integrate only over
one canonically ordered chamber and omit the factorial.
Feature collisions can reduce the number of distinct labelings, but those
collisions form a null set under product Lebesgue measure.  We write
\begin{equation}
  \mu=\sum_{n=0}^{N_{\max}}\mu_n,
\end{equation}
with $\mu_0$ a unit point mass at $\emptyconfig$.  For later use, write
$P_t=\mathcal L(X_t)$ for the law of $X_t$.  Whenever the restriction of
$P_t$ to $\state_n$ is absolutely continuous with respect to $\mu_n$, define
the stratum density by
\begin{equation}
  p_{t,n}
  =\frac{\dd(P_t|_{\state_n})}{\dd\mu_n},
  \qquad
  P_t(\dd X)=p_{t,n}(X)\mu_n(\dd X),
  \quad X\in\state_n.
  \label{eq:stratum-density}
\end{equation}
It integrates to the probability of the corresponding cardinality,
\begin{equation}
  \int_{\state_n}p_{t,n}(X)\mu_n(\dd X)
  =P_t(\state_n)=\Pp(N_t=n),
  \label{eq:stratum-density-mass}
\end{equation}
and therefore need not integrate to one.

The element $\dd z_{1:n}$ in \eqref{eq:reference-measure} is a volume
element, not a probability distribution.  A probability law is obtained
only after multiplying $\mu_n$ by the joint density $p_{t,n}$; that density
may couple all masses and features.  During reservoir dequantization the
observed $z_{1:n}$ are copied unchanged, so their Jacobian block is the
identity.  We integrate out the auxiliary $\xi$ in the raw likelihood, but
we do not integrate out $z_{1:n}$ because the features are observed.  A
weights-only marginal would additionally integrate over $\mathcal Z^n$
with respect to $\lambda_q^{\otimes n}$.

For any measurable $A\subseteq\state$, the complete law is reconstructed
from its stratum densities as
\begin{equation}
  P_t(A)=\sum_{n=0}^{N_{\max}}
  \int_{A\cap\state_n}p_{t,n}(X)\mu_n(\dd X).
  \label{eq:full-law-from-strata}
\end{equation}
Thus $\mu_n$ is a sigma-finite coordinate-volume measure, not a uniform
probability law on configurations.

The $1/n!$ convention is not cosmetic.  It determines the factors in the birth and merge rates derived in Section~\ref{sec:reverse}.

\subsection{A notation table}

\begin{center}
\small
\begin{tabular}{>{\raggedright\arraybackslash}p{0.20\linewidth}p{0.70\linewidth}}
\toprule
Symbol & Meaning \\
\midrule
$M$ & conserved total mass, fixed or conditioned upon \\
$\mathcal Z,\lambda_q$ & continuous feature space $\R^q$ and its Lebesgue reference measure \\
$N_{\max}$ & maximum allowed number of components \\
$v$ & mass in the reservoir \\
$(m_i,z_i)$ & mass and continuous feature of component $i$ \\
$N(X),N_t$ & number of components in state $X$, and $N(X_t)$ at time $t$ \\
$\eta_X$ & finite measure representing one weighted configuration \\
$\state_n,\state$ & $n$-component state space and its disjoint union \\
$\mu_n$ & unordered reference measure with factor $1/n!$ \\
$\nu_n^{\mathrm{raw}}$ & reference measure for unordered normalized raw data \\
$\Delta_{n-1}^{\circ}$ & open simplex of $n$ positive normalized weights \\
$q_\epsilon$ & fixed density of the auxiliary reservoir fraction $\xi$ \\
$P_t,p_{t,n}$ & forward marginal law and its density on $\state_n$ \\
$P_\theta,p_{\theta,n}$ & learned endpoint law and its density on $\state_n$ \\
$y$ & optional conditioning information \\
$T,\tau,S_N$ & terminal time, switch time, and nucleation duration $S_N=T-\tau$ \\
$\beta,\gamma,\delta$ & prescribed split, reservoir-flow, and deletion schedules \\
$\kappa_t$ & prescribed normalized split-mark density \\
$\splitmap_i,\mergemap_{ab}$ & split of component $i$ and merge of pair $a,b$ \\
$\deathmap_i,\birthmap_{m,z}$ & deletion and insertion maps \\
$Q_t,\overline Q_s$ & forward and reverse jump kernels \\
$\nu_X,\varrho$ & reference measures for forward histories and raw reverse records \\
$Y_s,\widetilde X_s$ & exact reversed state and reversed training path at reverse time $s$ \\
$L,\Pi_\theta$ & Stage-N leaf budget and its learned probability mass function \\
$A_\theta^{(\omega)},A_\theta^{(m)}$ & birth-mark densities in mass-fraction and absolute-mass coordinates \\
$B_L$ & fixed conditional birth hazard determined by $L$ \\
$\Lambda_\theta^C$ & learned total coagulation rate \\
$K_\theta(a,b\mid X,s)$ & learned merge intensity for unordered pair $a<b$ \\
$a_t(x)$ & optional control field at spatial coordinate $x=\operatorname{pos}(z)$ \\
$\lambda_i(t,X,a_t)$ & controlled evaporation rate of component $i$ \\
\bottomrule
\end{tabular}
\end{center}

\section{Conservative event maps}
\label{sec:event-maps}

This section defines the four operations used by the model.  Each operation preserves $\mathcal M(X)=M$ exactly.

\subsection{Mass exchange with the reservoir}

Between jumps, component features stay fixed and masses obey
\begin{align}
  \frac{\dd m_i}{\dd t}
  &=-\gamma(t)\frac{v}{M}m_i,
  \label{eq:forward-flow-mi}\\
  \frac{\dd v}{\dd t}
  &=\gamma(t)\frac{v}{M}\sum_{i=1}^{N(X)}m_i,
  \label{eq:forward-flow-v}
\end{align}
where $\gamma(t)\geq0$ is prescribed.  This is the forward evaporation flow.  It is deliberately simple: it transfers mass to the reservoir without changing relative component masses.

Let $c=M-v=\sum_i m_i$ and $r=c/(M-c)=c/v$.  On an interval with no jump,
\begin{equation}
  r(t)=r(t_0)\exp\!\left(-\int_{t_0}^{t}\gamma(u)\dd u\right),
  \qquad
  c(t)=M\frac{r(t)}{1+r(t)},
  \label{eq:closed-form-flow}
\end{equation}
and every $m_i$ is multiplied by $c(t)/c(t_0)$.  Thus the flow is available in closed form.
To verify the formula, sum \eqref{eq:forward-flow-mi} to obtain
\begin{equation}
  \frac{\dd c}{\dd t}=-\gamma(t)\frac{c(M-c)}{M}.
\end{equation}
For $0<c<M$, differentiating
$\log r=\log c-\log(M-c)$ gives
$\dd\log r/\dd t=-\gamma(t)$, which integrates to
\eqref{eq:closed-form-flow}.  The two boundary cases are fixed points:
$c=0$ is the empty state and $v=M$, while $v=0$ means $c=M$.  The odds
$r$ need not be evaluated at either boundary.

Running time backward changes the signs:
\begin{align}
  \frac{\dd m_i}{\dd s}
  &=+\gamma(T-s)\frac{v}{M}m_i,
  &
  \frac{\dd v}{\dd s}
  &=-\gamma(T-s)\frac{v}{M}\sum_i m_i.
  \label{eq:reverse-growth-flow}
\end{align}
Equation~\eqref{eq:reverse-growth-flow} is the growth stage of generation.  Because its vector field vanishes at $m_i=0$ and $v=0$, neither component mass nor reservoir mass can cross below zero.

\subsection{Binary fragmentation and coagulation}

Choose a component $(m,z)$, a fraction $\rho\in(0,1)$, and a displacement $u\in\R^q\setminus\{0\}$.  Define
\begin{align}
  m_a&=\rho m,
  &z_a&=z-(1-\rho)u,
  \label{eq:split-a}\\
  m_b&=(1-\rho)m,
  &z_b&=z+\rho u.
  \label{eq:split-b}
\end{align}
The split map $\splitmap_i(X;\rho,u)$ replaces component $i$ by these two children.  It preserves both mass and the mass-weighted feature centroid:
\begin{equation}
  m_a+m_b=m,
  \qquad
  m_a z_a+m_b z_b=mz.
  \label{eq:centroid-conservation}
\end{equation}

The inverse map merges an unordered pair $a<b$:
\begin{equation}
  \mergemap_{ab}(X)
  =X\setminus\{(m_a,z_a),(m_b,z_b)\}
   \cup\left\{\left(m_a+m_b,
   \frac{m_a z_a+m_b z_b}{m_a+m_b}\right)\right\}.
  \label{eq:merge-map}
\end{equation}

To avoid counting the same split twice, let $\mathcal U_+$ contain exactly one of $u$ and $-u$ for almost every nonzero $u$.  One choice is the set in which the first nonzero coordinate is positive.  We take split marks in
\begin{equation}
  \mathcal A=(0,1)\times\mathcal U_+.
\end{equation}
Given two distinct children, canonical orientation determines $u=z_b-z_a\in\mathcal U_+$ and then
\begin{equation}
  m=m_a+m_b,\qquad
  \rho=\frac{m_a}{m},\qquad
  z=\frac{m_a z_a+m_b z_b}{m}.
  \label{eq:inverse-split-coordinates}
\end{equation}

\subsection{Deletion and nucleation}

The deletion map transfers an entire component to the reservoir:
\begin{equation}
  \deathmap_i(X)
  =\left(v+m_i;\{(m_j,z_j):j\neq i\}\right).
  \label{eq:death-map}
\end{equation}
Its inverse is a birth or nucleation event.  For $0<m\leq v$,
\begin{equation}
  \birthmap_{m,z}(X)
  =\left(v-m;\{(m_i,z_i)\}_{i=1}^{N(X)}\cup\{(m,z)\}\right).
  \label{eq:birth-map}
\end{equation}
The equality case $m=v$ maps to the boundary $v=0$.  All coordinate-density
formulas below use the interior case $0<m<v$.  With reservoir-dequantized
data, finite flow times, and split fractions in $(0,1)$, the forward law
stays in that interior before deletion; reversing a deletion then indeed has
$m<v$ because the pre-deletion reservoir was strictly positive.
Unlike a model restricted to single-component insertion without a
pair-merge operation, births here are followed by learned pairwise merges.
The representation of a final component can therefore be assembled from
several nucleated pieces.

\begin{proposition}[Exact conservation of every operation]
\label{prop:event-conservation}
The flows in \eqref{eq:forward-flow-mi}--\eqref{eq:forward-flow-v} and
\eqref{eq:reverse-growth-flow}, and the maps
$\splitmap_i$, $\mergemap_{ab}$, $\deathmap_i$, and $\birthmap_{m,z}$ all preserve $\mathcal M(X)=M$.
\end{proposition}

\begin{proof}
For either flow, summing the derivatives of $v$ and all $m_i$ gives zero.  A split replaces $m$ by $\rho m+(1-\rho)m=m$, and a merge performs the reverse operation.  A deletion adds $m_i$ to $v$ before removing $m_i$; a birth subtracts the same amount from $v$ before inserting it.  Therefore every operation leaves the total unchanged.
\end{proof}

\subsection{A complete numerical example}

Take $M=10$ and $q=1$, so each $z_i$ is a scalar position.  Consider the
initial state
\begin{equation}
  X_0=\left(2;\{(3,1),(5,4)\}\right)\in\state_2.
  \label{eq:worked-example-initial}
\end{equation}
Here the reservoir has mass $2$, the two components have masses $3$ and $5$,
and their positions are $1$ and $4$.  The total is $2+3+5=10=M$.

Split the component $(m,z)=(5,4)$ using $\rho=0.4$ and $u=2$.  Equations
\eqref{eq:split-a}--\eqref{eq:split-b} give
\begin{equation}
  (5,4)\longmapsto(2,2.8),(3,4.8).
  \label{eq:worked-example-split}
\end{equation}
The resulting state is
\begin{equation}
  X_{t_1}=\left(2;\{(3,1),(2,2.8),(3,4.8)\}\right)
  \in\state_3.
  \label{eq:worked-example-after-split}
\end{equation}
The split preserves mass, and it also preserves the weighted feature because
$2(2.8)+3(4.8)=5(4)=20$.

Suppose a subsequent flow segment reduces the total component mass from
$8$ to $6$.  The basic flow preserves relative masses, so it multiplies
each component mass by $6/8=3/4$ and transfers the remaining mass to the
reservoir:
\begin{equation}
  X_{t_2}
  =\left(4;\{(2.25,1),(1.5,2.8),(2.25,4.8)\}\right)
  \in\state_3.
  \label{eq:worked-example-after-flow}
\end{equation}
The continuous flow changes $v$ and the masses but not the cardinality or
features.  If the component $(1.5,2.8)$ is then deleted, its mass enters the
reservoir and
\begin{equation}
  X_{t_3}
  =\left(5.5;\{(2.25,1),(2.25,4.8)\}\right)
  \in\state_2.
  \label{eq:worked-example-after-deletion}
\end{equation}
Every displayed state has total mass $10$.  Reversing these operations first
births $(1.5,2.8)$, then reverses the flow, and finally merges
$(2,2.8)$ with $(3,4.8)$.  The merged component has
\begin{equation}
  m=2+3=5,
  \qquad
  z=\frac{2(2.8)+3(4.8)}{5}=4,
\end{equation}
so the reverse sequence recovers \eqref{eq:worked-example-initial} exactly.

The continuous flow and four jump operations move between strata as follows:
\begin{center}
\begin{tabular}{lll}
\toprule
Operation & Stratum change & Quantities changed directly \\
\midrule
Continuous flow & $\state_n\to\state_n$ & $v,m_i$ \\
Split & $\state_n\to\state_{n+1}$ & one component becomes two \\
Deletion & $\state_n\to\state_{n-1}$ & one component moves into $v$ \\
Birth & $\state_n\to\state_{n+1}$ & mass moves out of $v$ \\
Merge & $\state_n\to\state_{n-1}$ & two components become one \\
\bottomrule
\end{tabular}
\end{center}

\section{The forward fragmentation--evaporation process}
\label{sec:forward}

\subsection{Two stages}

Choose $0<\tau<T$.  The forward process starts at $X_0\sim\pdata$ and has two stages.

\paragraph{Stage F: fragmentation, $0\leq t\leq\tau$.}
Each component splits at rate $\beta(t)$, where
$0\leq\beta(t)\leq\beta_{\max}<\infty$, provided that
$N(X)<N_{\max}$.  Conditional on a split, the parent is therefore uniform
among the current components.  The split mark has a known normalized density
\begin{equation}
  \kappa_t(\rho,u\mid X,i)
  \quad\text{on }\mathcal A.
  \label{eq:split-kernel}
\end{equation}
A ``density on $\mathcal A$'' means a nonnegative function integrating to
one relative to $\dd\rho\,\lambda_q(\dd u)$ restricted to
$\mathcal U_+$.  It may depend jointly on the full state and the selected
parent; no factorization of $\rho$ and $u$ is assumed.
A simple choice is a symmetric Beta density for $\rho$ and a half-space Gaussian density for $u$.  The evaporation flow \eqref{eq:forward-flow-mi}--\eqref{eq:forward-flow-v} acts between events.  There are no deletion events in this stage.

\paragraph{Stage E: evaporation and deletion, $\tau<t<T$.}
Splitting is turned off.  Every remaining component is deleted with hazard
\begin{equation}
  \delta(t)=\frac{1}{T-t}.
  \label{eq:death-hazard}
\end{equation}
The conservative evaporation flow continues between deletions.  Because
\begin{equation}
  \exp\!\left(-\int_\tau^t\delta(u)\dd u\right)
  =\frac{T-t}{T-\tau},
  \label{eq:uniform-death-survival}
\end{equation}
the deletion time of each component alive at $\tau$ is uniform on $(\tau,T)$.  This observation gives an efficient exact simulator.

\subsection{Jump generator}

For a bounded test function $f:\state\to\R$, the jump part of the forward generator is
\begin{align}
  (\mathcal J_t f)(X)
  ={}&\ind\{t\leq\tau,\,N(X)<N_{\max}\}
  \sum_{i=1}^{N(X)}\beta(t)
  \int_{\mathcal A}
  \left[f\!\left(\splitmap_i(X;\rho,u)\right)-f(X)\right]
  \kappa_t(\rho,u\mid X,i)\dd\rho\dd u
  \nonumber\\
  &+\ind\{t>\tau\}
  \sum_{i=1}^{N(X)}\delta(t)
  \left[f\!\left(\deathmap_i(X)\right)-f(X)\right].
  \label{eq:forward-jump-generator}
\end{align}
The full generator adds the directional derivative along the deterministic mass flow.  This makes $X_t$ a time-inhomogeneous piecewise-deterministic Markov process.

\begin{theorem}[Well-posedness and exact terminal state]
\label{thm:forward-terminal}
Suppose $N(X_0)\leq N_{\max}$ almost surely,
$\beta(t)\leq\beta_{\max}<\infty$, and
$\int_0^T\gamma(t)\dd t<\infty$.  Then:
\begin{enumerate}
  \item the forward process has finitely many jumps on every interval $[0,t]$ with $t<T$;
  \item masses and reservoir mass remain nonnegative and sum to $M$;
  \item $X_t\to\emptyconfig$ almost surely as $t\uparrow T$.  We may therefore set $X_T=\emptyconfig$.
\end{enumerate}
\end{theorem}

\begin{proof}
During Stage F, a state with $n$ components has total split rate at most
$n\beta_{\max}\leq N_{\max}\beta_{\max}$.  Moreover, each split raises
cardinality by one and splitting stops at $N_{\max}$.  Thus there are at most
$N_{\max}-N(X_0)$ splits.

At the boundary $m_i=0$, the derivative in \eqref{eq:forward-flow-mi} is zero.  At $v=0$, both derivatives in \eqref{eq:forward-flow-mi}--\eqref{eq:forward-flow-v} are zero.  Hence the flow cannot leave the nonnegative simplex.  Proposition~\ref{prop:event-conservation} shows that jumps also preserve the simplex and total mass.

There are finitely many components at time $\tau$.  By \eqref{eq:uniform-death-survival}, each one has probability zero of surviving until $T$.  A finite union of probability-zero survival events still has probability zero.  Consequently all components are deleted before $T$ almost surely, and every deletion transfers its remaining mass to $v$.  The terminal state is therefore $(M;\varnothing)$.
\end{proof}

\begin{remark}[Why the singular-looking hazard is useful]
The hazard $1/(T-t)$ diverges, but every realized deletion time is strictly smaller than $T$ almost surely.  The process has only finitely many components in Stage E, so it still has finitely many deletion events.  The divergence solves a real modeling problem: a bounded deletion hazard would leave a nonzero probability of a nonempty terminal state.
\end{remark}

\section{Finite-horizon time reversal}
\label{sec:reverse}

\subsection{The measure-level formula}

The safest reversal formula is written using measures, before choosing coordinates.

\begin{definition}[Forward jump flux]
Let $Q_t(X,\dd Y)$ be the forward jump-rate kernel.  The jump flux at time $t$ is the measure
\begin{equation}
  \Gamma_t(\dd X,\dd Y)
  =P_t(\dd X)Q_t(X,\dd Y).
  \label{eq:forward-flux}
\end{equation}
It describes probability per unit time moving from $X$ to $Y$.
\end{definition}

\begin{assumption}[Regularity for finite-horizon reversal]
\label{ass:reversal-regularity}
The state space is standard Borel.  The forward process is a Borel-right,
non-explosive piecewise-deterministic Markov process.  Within each stratum,
the flow before it reaches a boundary is a measurable $C^1$ bijection.
Interior jumps are not simultaneous and do not occur at predictable times.
The transposed incoming jump flux is absolutely continuous with respect to
the current marginal law, and the marginal laws and fluxes are regular in
time wherever the rates below are evaluated.
\end{assumption}

\begin{theorem}[Finite-horizon reverse jump kernel]
\label{thm:reverse-kernel}
Under Assumption~\ref{ass:reversal-regularity}, define
$\overline Q_{T-t}$ by the disintegration
\begin{equation}
  P_t(\dd Y)\overline Q_{T-t}(Y,\dd X)
  =P_t(\dd X)Q_t(X,\dd Y).
  \label{eq:flux-reversal}
\end{equation}
Set $Y_0=X_T$, $Y_T=X_0$, and, for $0<s<T$, use the c\`adl\`ag
reversal
\begin{equation}
  Y_s=X_{(T-s)-}.
  \label{eq:cadlag-reversal}
\end{equation}
Then $Y$ is Markov.  Between jumps it follows the inverse deterministic
flow, and its jump-rate kernel is $\overline Q_s$.
\end{theorem}

\begin{proof}[Proof sketch]
Fix $t$ and a small $h>0$.  To first order in $h$, the joint measure of a forward jump from $X$ at time $t-h$ to $Y$ at time $t$ is
\[
  h\,P_{t-h}(\dd X)Q_{t-h}(X,\dd Y)+o(h).
\]
Under the stated time regularity, replacing both $P_{t-h}$ and $Q_{t-h}$
by their time-$t$ versions changes only the $o(h)$ term.  Bayes' rule
conditions this joint measure on the endpoint $Y$.  The resulting
conditional rate is exactly the disintegration in
\eqref{eq:flux-reversal}.  When no jump occurs, a forward deterministic
segment is traversed backward, so its velocity changes sign.  Applying the
same argument at each time proves the Markov property and identifies the
reverse characteristics.
The stated regularity justifies the disintegration and passage from the
short-time identity to the Markov characteristics; see
\cite{conforti2022time} for a general treatment.  The left limit in
\eqref{eq:cadlag-reversal} makes the reversed path right-continuous: at a
forward jump time $t$, the reverse pre-jump state is $X_t$ and the reverse
post-jump state is $X_{t-}$.
\end{proof}

Equation~\eqref{eq:flux-reversal} is a finite-time statement.  It depends on the current marginal $P_t$.  It is not the same as equilibrium detailed balance.

\begin{remark}[The terminal point mass is an entrance law]
At $t=T$, the forward law is the point mass at $\emptyconfig$, so
stratum-wise density ratios are used only for $0<t<T$.  The exact reverse at
$s=0$ is understood as the entrance-law limit of $X_{(T-s)-}$ as
$s\downarrow0$.
This limit is unambiguous even though the coordinate density formula is not
defined at the point mass itself.  A numerical implementation may instead
start at reverse time $s=\varepsilon$ from $\emptyconfig$, where
$0<\varepsilon<T-\tau$.  By
\eqref{eq:uniform-death-survival} and a union bound,
\begin{equation}
  \TV\!\left(P_{T-\varepsilon},\delta_{\emptyconfig}\right)
  =\Pp(N_{T-\varepsilon}>0)
  \leq
  \frac{\E[N(X_\tau)]\,\varepsilon}{T-\tau}
  \leq\frac{N_{\max}\varepsilon}{T-\tau}.
  \label{eq:entrance-tv-bound}
\end{equation}
Markov kernels contract total variation, so the same quantity bounds the
endpoint error caused solely by this $\varepsilon$-start approximation when
the remaining reverse dynamics are exact.
\end{remark}

\subsection{Coordinate Jacobian for a split}

For one parent, the coordinate map is
\begin{equation}
  (m,\rho,z,u)\longmapsto(m_a,m_b,z_a,z_b)
\end{equation}
with definitions \eqref{eq:split-a}--\eqref{eq:split-b}.

\begin{lemma}[Split Jacobian]
\label{lem:split-jacobian}
The absolute determinant of the above map is
\begin{equation}
  \left|\frac{\partial(m_a,m_b,z_a,z_b)}
  {\partial(m,\rho,z,u)}\right|=m.
  \label{eq:split-jacobian}
\end{equation}
\end{lemma}

\begin{proof}
The mass block $(m,\rho)\mapsto(\rho m,(1-\rho)m)$ has absolute determinant $m$.  For each feature coordinate, the map $(z,u)\mapsto(z-(1-\rho)u,z+\rho u)$ has matrix
\[
  \begin{pmatrix}1&-(1-\rho)\\1&\rho\end{pmatrix}
\]
and determinant one.  The full derivative is block lower triangular because the child masses do not depend on $(z,u)$.  Multiplying the determinants gives $m$.
\end{proof}

\begin{assumption}[Densities for coordinate formulas]
\label{ass:stratum-densities}
For every relevant $n$ and $0<t<T$, the marginal restricted to a stratum is
dominated by its reference measure:
\begin{equation}
  P_t|_{\state_n}\ll\mu_n.
\end{equation}
Its density $p_{t,n}$ is positive at every state where a density ratio is
evaluated.  The measure identity in Theorem~\ref{thm:reverse-kernel} does not
need this assumption; only the coordinate ratios below do.  The random
reservoir augmentation in \eqref{eq:reservoir-dequantization} is one way to
avoid a singular data law.
\end{assumption}

The assumption rules out atoms on $v=0$ and singular feature laws.  It is
preserved by the intended interior construction under ordinary smoothness
conditions: for a finite flow interval,
$m_i(t)=m_i(t_0)\exp[-\int\gamma(u)v(u)/M\,\dd u]>0$; an initially positive
reservoir remains positive; split fractions lie in $(0,1)$; and continuous
split and birth marks hit $u=0$ or exact feature collisions only on null
sets.  Absolute continuity nevertheless remains an explicit assumption,
because arbitrary data or mark kernels need not satisfy it.

\subsection{Explicit birth and merge rates}

Let $s=T-t$ denote reverse time.  We now write the exact reverse rates under the reference convention \eqref{eq:reference-measure}.

\begin{proposition}[Exact reverse of deletion]
\label{prop:reverse-death}
Let $Y\in\state_{n-1}$ and let $X=\birthmap_{m,z}(Y)\in\state_n$.  Define
$d_t(X;m,z)$ as the forward per-component deletion intensity of the
component $(m,z)$ in $X$; in the present model it is
$\ind\{t>\tau\}\delta(t)$.  Then the reverse birth intensity, as a density
with respect to $\dd m\,\lambda_q(\dd z)$, is
\begin{equation}
  a_s^*(m,z\mid Y)
  =d_t(X;m,z)\frac{p_{t,n}(X)}{p_{t,n-1}(Y)},
  \qquad t=T-s.
  \label{eq:exact-birth-rate}
\end{equation}
There is no extra factor of $n$: the sum over the $n$ possible deleted labels cancels the ratio between $1/n!$ and $1/(n-1)!$.
\end{proposition}

\begin{proof}
Consider all labeled representations of $X$ and sum the forward deletion flux over the deleted index.  Symmetry makes the $n$ terms equal, while the reference measure contributes $1/n!$.  Thus $n/n!=1/(n-1)!$, which is precisely the factor in the reference measure of $Y$.  The coordinate transformation $(Y,m,z)\mapsto X$ has determinant one.  Substituting these facts into \eqref{eq:flux-reversal} gives \eqref{eq:exact-birth-rate}.
\end{proof}

\begin{proposition}[Exact reverse of fragmentation]
\label{prop:reverse-split}
Let $X\in\state_n$ with $n<N_{\max}$ and let
$Y=\splitmap_i(X;\rho,u)\in\state_{n+1}$.  Suppose the forward marked split intensity is
\begin{equation}
  r_t^{\mathrm{split}}(X,i,\rho,u)
  =\beta(t)\kappa_t(\rho,u\mid X,i)
  \label{eq:marked-split-intensity}
\end{equation}
with canonical $u\in\mathcal U_+$.  Let $a<b$ be the resulting pair in $Y$ and let $m=m_a+m_b$.  Then the exact reverse merge intensity for this unordered pair is
\begin{equation}
  k_s^*(a,b\mid Y)
  =\frac{1}{m}\,
  r_t^{\mathrm{split}}(X,i,\rho,u)
  \frac{p_{t,n}(X)}{p_{t,n+1}(Y)},
  \qquad t=T-s.
  \label{eq:exact-merge-rate}
\end{equation}
No additional cardinality or factor-of-two correction appears under the
canonical-orientation convention.  The formula is stated on the
full-measure set $z_a\neq z_b$; the excluded diagonal corresponds to
$u=0$.
\end{proposition}

\begin{proof}
After selecting the parent, the forward reference factor is
$n/n!=1/(n-1)!$.  If child coordinates are first integrated as an ordered
pair, selecting an unordered pair contributes
$\binom{n+1}{2}/(n+1)!=1/[2(n-1)!]$.  However, an unordered pair has two
ordered representations in the labeled-coordinate integral.  That full
integral is twice its representation on the canonical half-space
$u\in\mathcal U_+$, which cancels the factor $1/2$.  Thus both the
selected-parent measure and the canonically oriented
selected-pair measure carry $1/(n-1)!$.  By
Lemma~\ref{lem:split-jacobian},
$\dd m_a\dd m_b\dd z_a\dd z_b=m\,\dd m\dd\rho\dd z\dd u$.
Equating the two flux densities in \eqref{eq:flux-reversal} therefore
multiplies the forward marked intensity by $1/m$.  The remaining factor is
the marginal density ratio $p_{t,n}(X)/p_{t,n+1}(Y)$.
\end{proof}

The density ratios in \eqref{eq:exact-birth-rate} and
\eqref{eq:exact-merge-rate} are not available for real data.
Section~\ref{sec:learning} derives a simulation-based objective whose
population optimum recovers the same reverse event law under the joint
realizability and support assumptions of
Theorem~\ref{thm:intensity-consistency}, without evaluating these ratios.

\subsection{Why detailed balance is not the derivation}

For a time-homogeneous process with stationary law $\pi$, detailed balance is
the equality of measures
\begin{equation}
  \pi(\dd X)Q(X,\dd Y)=\pi(\dd Y)Q(Y,\dd X).
  \label{eq:detailed-balance}
\end{equation}
For example, applying the same split coordinates as above would give
\begin{equation}
  k^{\mathrm{merge}}(Y,a,b)
  =\frac{1}{m_a+m_b}\,
    r^{\mathrm{split}}(X,\rho,u)
    \frac{\pi_n(X)}{\pi_{n+1}(Y)}.
  \label{eq:stationary-merge-rate}
\end{equation}
After expressing mass and feature coordinates in fixed reference units, if
$\pi_n(X)\propto\exp[-\beta_{\mathrm{th}} F_n(X)]$ relative to the same stratum
measures, then
\begin{equation}
  \log\frac{k^{\mathrm{merge}}(Y,a,b)}
                 {r^{\mathrm{split}}(X,\rho,u)}
  =\beta_{\mathrm{th}}\{F_{n+1}(Y)-F_n(X)\}
   -\log\!\left(\frac{m_a+m_b}{m_0}\right),
  \label{eq:gibbs-rate-ratio}
\end{equation}
where $m_0>0$ is a reference mass; its constant contribution can be
absorbed into the rate convention.  This expression is still subject to
any explicitly chosen base-measure terms.  Our process is instead
time-inhomogeneous and transports a data distribution to an absorbing state.
Its reverse depends on $p_{t,n}$, not on a stationary Gibbs law.  A learned
free energy may still regularize rates, but it cannot replace the
finite-horizon flux calculation.

\section{Learning from reversed forward paths}
\label{sec:learning}

\subsection{Latent histories}

A forward history $H$ contains
\begin{itemize}
  \item all split times, selected parents, and split marks $(\rho,u)$;
  \item all deletion times and the deleted component states just before deletion;
  \item the deterministic flow segments between events;
  \item persistent lineage identifiers, used only to identify inverse merge pairs.
\end{itemize}
Conditional on $(X,$ event times, split marks, and discrete choices$)$, the
recorded pre-deletion component states and flow segments are deterministic.
They are stored for implementation convenience but are not additional
coordinates of the history reference measure.
By Theorem~\ref{thm:forward-terminal}, every history ends at $\emptyconfig$.  Reverse the event order and replace every deletion by a birth and every split by a merge.  The resulting history is denoted $R(H)$.  It begins at $\emptyconfig$ and ends at the original data state.

Let $Q_X^F$ be the forward-history law conditional on $X_0=X$, with density
$q_X^F(H)$ relative to a history reference measure $\nu_X(\dd H)$.  It is
important to reverse the \emph{joint} coordinates $(X,H)$, not $H$ alone.
More precisely, the history space is a disjoint union over the finite
numbers of events and their discrete lineage patterns.  On each piece,
$\nu_X$ is Lebesgue measure on the ordered event-time simplices and
continuous split marks, multiplied by counting measure on parent, deletion,
and lineage choices.  The raw reverse-record space is defined in the same
way: its reference measure $\varrho(\dd R)$ is Lebesgue measure on ordered
birth and merge times and continuous birth marks, multiplied by counting
measure on the event counts and merge-pair choices.

Choose a fixed Borel ordering of component coordinates, such as
lexicographic order, away from its measure-zero set of ties.  On this
canonical chamber, the quotient measure $\mu_n$ is ordinary coordinate
volume: the $n!$ labeled chambers in
\eqref{eq:reference-measure} cancel its factor $1/n!$.  New lineages are
assigned by event order and canonical child orientation.  Arbitrary lineage
renamings are therefore not additional outcomes.  The factor $1/n!$ occurs
only in $\mu_n$.  Likewise, $\varrho$ is plain volume on the ordered
birth-time simplex; the factor $L!$ below is part of its probability density
and occurs only once.
Define
\begin{equation}
  \Psi:(X,H)\longmapsto R,
  \label{eq:joint-reversal-map}
\end{equation}
where $R$ is the raw reverse event record, including event times, birth
marks, and merge choices.  The endpoint coordinates in $X$ become the birth
marks missing from a conditional forward history.

This resolves a simple dimension check.  Suppose $X$ has $n$ components,
the forward history has $C$ splits, and therefore $L=n+C$ deletions.  With
discrete lineages fixed, the continuous dimension on the forward side is
\begin{equation}
  n(q+1)+C(q+2)+L,
\end{equation}
while the reverse record has dimension
\begin{equation}
  L(q+2)+C.
\end{equation}
These are equal because $L=n+C$.  On each fixed lineage pattern, $\Psi$ is
a square, almost-everywhere differentiable bijection.  Its absolute
Jacobian is denoted $J_\Psi(X,H)$.  Appendix~\ref{app:forward-density}
gives an explicit product formula.

\begin{lemma}[Joint history change of variables]
\label{lem:joint-history-change}
Outside simultaneous events, feature collisions, and coordinate-ordering
ties, the map $\Psi$ is a bijection from admissible endpoint--forward-history
pairs to admissible raw reverse records.  If
$\mathcal L_\theta$ is a normalized density relative to $\varrho$, then
\begin{equation}
  \sum_{n=0}^{N_{\max}}\int_{\state_n}\int
  \mathcal L_\theta(\Psi(X,H))J_\Psi(X,H)
  \nu_X(\dd H)\mu_n(\dd X)=1.
  \label{eq:joint-change-normalization}
\end{equation}
\end{lemma}

\begin{proof}
Given $(X,H)$, reverse the event order, turn each deletion into a birth, and
turn each split into a merge.  This constructs one record $R$.  Conversely,
simulate any admissible $R$ from $\emptyconfig$ to obtain its endpoint $X$;
then reverse every deterministic segment and event to recover $H$.  The
canonical component ordering and lineage rules make these constructions
inverse to each other.  The excluded cases are null under the continuous
reference measures.

On a fixed lineage piece, ordinary multivariate change of variables gives
the integrand in \eqref{eq:joint-change-normalization}; the dimensions match
by the calculation above.  Summing first over its discrete patterns and then
over $n$ gives
$\int\mathcal L_\theta(R)\varrho(\dd R)=1$.  Because $\mu_n$ equals
coordinate volume on the canonical chamber, no extra $n!$ appears.  Because
event times are already ordered in $\varrho$, no second $L!$ appears.
\end{proof}

\subsection{Path-space evidence lower bound}

\begin{theorem}[Path-space ELBO]
\label{thm:path-elbo}
Let $\mathcal L_\theta(R)$ be the density of a raw reverse event record
relative to $\varrho$.  For $X\in\state_n$, define the joint
endpoint--history density, relative to
$\mu_n(\dd X)\nu_X(\dd H)$, by
\begin{equation}
  p_{\theta,n}(X,H)
  =\mathcal L_\theta(\Psi(X,H))J_\Psi(X,H).
  \label{eq:joint-generative-density}
\end{equation}
Let
$p_{\theta,n}(X)=\int p_{\theta,n}(X,H)\nu_X(\dd H)$, including the sum
over discrete lineage patterns.  Lemma~\ref{lem:joint-history-change} shows
that these are the stratum densities of a normalized endpoint law.  Then,
whenever the supports are compatible,
\begin{equation}
  \log p_{\theta,n}(X)
  \geq
  \E_{H\sim Q_X^F}
  \left[
    \log\mathcal L_\theta(\Psi(X,H))
    +\log J_\Psi(X,H)
    -\log q_X^F(H)
  \right].
  \label{eq:path-elbo}
\end{equation}
Equality holds when $Q_X^F$ is the model posterior over $H$ given $X$.
\end{theorem}

\begin{proof}
Since $q_X^F$ is normalized,
\begin{align*}
  p_{\theta,n}(X)
  &=\int p_{\theta,n}(X,H)\nu_X(\dd H)\\
  &=\E_{H\sim Q_X^F}
  \left[\frac{p_{\theta,n}(X,H)}{q_X^F(H)}\right].
\end{align*}
Insert \eqref{eq:joint-generative-density}, take logarithms, and apply
Jensen's inequality.  Equality requires the importance ratio to be constant,
which is exactly the posterior condition.
\end{proof}

In the core model with prescribed flow and corruption,
only $\log\mathcal L_\theta$ depends on $\theta$.  Thus training needs only
sampled forward histories and the trainable reverse-record likelihood.  A
numerical ELBO additionally uses $q_X^F$, $J_\Psi$, and, for normalized raw
data, the dequantization terms from \eqref{eq:reservoir-dequantization}.

\subsection{Learned reverse event model}

Reverse time has two stages in the opposite order.

\paragraph{Auxiliary leaf budget.}
Let
\begin{equation}
  L=N(X_\tau)\in\{0,\ldots,N_{\max}\}
  \label{eq:leaf-budget}
\end{equation}
be the number of components present when forward deletion begins.  The
generative model first samples
\begin{equation}
  L\sim\Pi_\theta(L\mid M,y),
  \label{eq:leaf-count-model}
\end{equation}
where $y$ denotes optional conditioning information.  This is an auxiliary
pre-coagulation count, not the final output cardinality.

\paragraph{Stage N: nucleation, $0<s<S_N=T-\tau$.}
Conditional on $L$, the reversed deletion times are the order statistics of
$L$ independent $\operatorname{Uniform}(0,S_N)$ variables.  Equivalently,
when $N(X)<L$, the next-birth hazard is
\begin{equation}
  B_L(X,s)=\frac{L-N(X)}{S_N-s},
  \qquad 0\leq s<S_N,
  \label{eq:fixed-birth-bridge}
\end{equation}
and it is zero when $N(X)=L$.  The apparent deadline singularity is benign:
there are exactly $L\leq N_{\max}$ births, all strictly before $S_N$ almost
surely.  The state follows the growth flow between them.  Conditional on a
birth, its mass and feature have learned normalized density
\begin{equation}
  A_\theta^{(m)}(m,z\mid X,s,L),
  \qquad 0<m<v.
  \label{eq:birth-mark-model}
\end{equation}
The timing rule is fixed; only the leaf-count law and birth marks are learned.
Marginalizing $L$ gives a valid latent-mixture birth process, but it need not
be Markov in $X$ alone because the posterior law of $L$ can depend on the
past.  Keeping $(X,L)$ as the augmented state makes the process Markov and
makes exact birth-count support and non-explosion explicit.

\paragraph{Stage C: coagulation, $T-\tau<s<T$.}
The growth flow continues.  Each unordered pair $a<b$ has merge intensity
\begin{equation}
  K_\theta(a,b\mid X,s)\geq0,
  \label{eq:merge-model}
\end{equation}
and the total event rate is
\begin{equation}
  K_\theta^{\mathrm{tot}}(X,s)
  =\sum_{a<b}K_\theta(a,b\mid X,s).
  \label{eq:total-merge-rate}
\end{equation}

For a reversed training path, use the c\`adl\`ag convention
$\widetilde X_s=X_{(T-s)-}$ in the interior, with the endpoint conventions
of \eqref{eq:cadlag-reversal}.  Let $\mathcal E_B$ and $\mathcal E_C$ be its
birth and merge events.  Its trainable negative log likelihood is
\begin{align}
  \mathcal L_{\mathrm{path}}(\theta;H)
  ={}&
  -\log\Pi_\theta(L\mid M,y)
  -\sum_{e\in\mathcal E_B}
    \log A_\theta^{(m)}(m_e,z_e\mid\widetilde X_{s_e^-},s_e,L)
  \nonumber\\
  &+\int_{S_N}^{T}K_\theta^{\mathrm{tot}}(\widetilde X_s,s)\dd s
  -\sum_{e\in\mathcal E_C}
  \log K_\theta(a_e,b_e\mid\widetilde X_{s_e^-},s_e).
  \label{eq:path-loss}
\end{align}
The ordered birth times have fixed density $L!/S_N^L$; its hazard and
survival terms do not depend on $\theta$ and are omitted from
\eqref{eq:path-loss}.  Flow Jacobians and prescribed forward terms are also
fixed in $\theta$ and appear only in the numerical ELBO.

The merge survival integral is essential.  Training only observed pair
labels would learn which merge occurred but not its timing or stopping law.

\subsection{Unbiased integral estimator}

For any interval $[a,b]$ and path-dependent rate $\Lambda_\theta(s)$,
\begin{equation}
  \int_a^b\Lambda_\theta(s)\dd s
  =(b-a)\E_{U\sim\mathrm{Unif}(a,b)}[\Lambda_\theta(U)].
  \label{eq:uniform-integral-estimator}
\end{equation}
Thus one or more uniform times give an unbiased estimator of the merge
survival integral.  If the merge stage is divided into event-free intervals
$I_j=[a_j,b_j]$, a one-sample stratified estimator is
\begin{equation}
  \sum_j |I_j|\,K_\theta^{\mathrm{tot}}(\widetilde X_{U_j},U_j),
  \qquad U_j\sim\operatorname{Uniform}(I_j).
  \label{eq:stratified-integral-estimator}
\end{equation}
The interval-length factors are required for unbiasedness.

\subsection{Population consistency}

The following theorem states the model-class condition behind the consistency
claim.

\begin{theorem}[Fisher consistency of factorized reverse learning]
\label{thm:intensity-consistency}
Let $\Pi^*$, $A^{*(m)}$, and $K^*$ be the exact leaf-count law, conditional
birth-mark densities, and merge intensities of the augmented reverse
process.  Assume \emph{joint realizability}: there is one parameter value
$\theta^*$ for which all three model factors equal these exact factors
simultaneously.  Also assume that the learned factors assign positive
density or intensity wherever the exact factors are positive, and that the
displayed cross-entropies and compensator integrals are finite.  Then the population risk of
\eqref{eq:path-loss} is minimized by
$\Pi_\theta=\Pi^*$, $A_\theta^{(m)}=A^{*(m)}$, and $K_\theta=K^*$ almost everywhere
under the exact reverse path law.
\end{theorem}

\begin{proof}
Subtract the risk at the exact factors.  The leaf-count term becomes the
expected categorical divergence
$\E[\KL(\Pi^*\,\|\,\Pi_\theta)]$.  Conditional on every birth pre-state,
time, and $L$, the birth-mark term becomes
$\KL(A^{*(m)}\,\|\,A_\theta^{(m)})$; summing over births and taking expectation keeps it
nonnegative.  For the merge process, the compensator identity turns the
excess risk into
\begin{equation}
  \E\int_{S_N}^{T}\sum_{a<b}
  \left[
    K_\theta-K^*+K^*\log\frac{K^*}{K_\theta}
  \right]\dd s.
  \label{eq:intensity-divergence}
\end{equation}
For positive $a,b$, $a-b+b\log(b/a)\geq0$, with equality only at $a=b$.
We use the standard conventions $0\log(0/a)=0$ and assign infinite
divergence when $b>0$ but $a=0$.  Every term is therefore nonnegative, and
all vanish exactly at the stated jointly realizable factors.
\end{proof}

The theorem is conditional on joint realizability; it does not claim that a
finite network recovers the exact factors from finite data.  In particular,
the bounded-rate implementation below may exclude an exact bridge whose
rate diverges near the endpoint.  The theorem does show why no
explicit estimate of the marginal density ratios in
\eqref{eq:exact-birth-rate}--\eqref{eq:exact-merge-rate} is required.

\section{Neural parameterization}
\label{sec:parameterization}

\subsection{Permutation-equivariant encoder}

Let $M_{\mathrm{ref}}>0$ be a fixed reference mass in the same units as
$M$, and write $\ell_M=\log(M/M_{\mathrm{ref}})$.  This gives the networks a
dimensionless total-mass input when examples have different $M$.
For every component, form an embedding
\begin{equation}
  h_i=\phi_\theta\!\left(z_i,\frac{m_i}{M},s,\ell_M,y\right),
  \label{eq:element-embedding}
\end{equation}
where $y$ denotes optional conditioning information.  A Deep Sets or Set Transformer backbone produces updated embeddings and a pooled summary
\begin{equation}
  (\widehat h_1,\ldots,\widehat h_n,g)
  =\mathrm{SetNet}_\theta(h_1,\ldots,h_n).
  \label{eq:set-backbone}
\end{equation}
The component outputs are permutation equivariant and $g$ is invariant
\cite{zaheer2017deepsets,lee2019settransformer}.
For the empty set, define $g=g_\varnothing$ using a learned null embedding
(or a fixed zero vector followed by a network that also receives
$v/M,s,\ell_M,L,y$).  This makes the first birth-mark distribution in
Algorithm~3 well defined.

\subsection{Leaf count and birth marks}

The leaf-count head is a categorical distribution on a finite support:
\begin{equation}
  \Pi_\theta(\ell\mid M,y)
  =\operatorname{softmax}\!\left(a_\theta(\ell_M,y)\right)_\ell,
  \qquad \ell=0,\ldots,N_{\max}.
  \label{eq:leaf-count-head}
\end{equation}
For theoretical statements, the categorical family includes boundary points
with exact zero probabilities.  In code, a masked softmax enforces known
structural zeros; an ordinary softmax represents the interior and approaches
an unknown zero only in the limit.

Write the birth mass as a fraction $\omega=m/v\in(0,1)$.  A fully
normalized, directly implementable mark head is a mixture
\begin{equation}
  A_\theta^{(\omega)}(\omega,z\mid X,s,L)
  =\sum_{j=1}^{J_{\mathrm{mix}}}\alpha_j^{\mathrm{mix}}
  \operatorname{Beta}(\omega;\alpha_j,\beta_j)
  \mathcal N(z;\bar z_j,\operatorname{diag}(\sigma_j^2)),
  \label{eq:birth-mixture}
\end{equation}
where all parameters depend on $(g,v/M,s,\ell_M,L,y)$, the mixture weights
$\alpha_j^{\mathrm{mix}}$ use a softmax, and positive parameters use
softplus.  The density with respect to $\dd m\dd z$ is
\begin{equation}
  A_\theta^{(m)}(m,z\mid X,s,L)
  =\frac{1}{v}A_\theta^{(\omega)}(m/v,z\mid X,s,L).
  \label{eq:birth-density-jacobian}
\end{equation}
For high-dimensional features, the Gaussian mixture may be replaced by a conditional normalizing flow, provided its density and samples are tractable.

\subsection{Merge intensity}

For every candidate unordered pair, compute a symmetric logit
\begin{equation}
  \ell_{ab}
  =\psi_\theta(\widehat h_a+\widehat h_b,
                |\widehat h_a-\widehat h_b|,
                g,s,\ell_M,y)
  =\ell_{ba}.
  \label{eq:pair-logit}
\end{equation}
When $N(X)\geq2$, let
\begin{equation}
  \pi_\theta(a,b\mid X,s)
  =\frac{\exp(\ell_{ab})}{\sum_{i<j}\exp(\ell_{ij})}
  \label{eq:pair-softmax}
\end{equation}
and parameterize a finite total rate with bounded average rate per pair,
\begin{equation}
  \Lambda_\theta^C(X,s)
  =k_{\max}(s)\binom{N(X)}{2}
  \sigma\!\left(g_\theta(g,s,\ell_M,y)\right).
  \label{eq:bounded-total-merge}
\end{equation}
When $N(X)<2$, set $\Lambda_\theta^C=0$ and do not evaluate the pair softmax.
Then
\begin{equation}
  K_\theta(a,b\mid X,s)
  =\Lambda_\theta^C(X,s)\pi_\theta(a,b\mid X,s).
  \label{eq:factorized-merge-rate}
\end{equation}
Because $\pi_\theta$ is normalized over the active candidate set,
$K_\theta^{\mathrm{tot}}=\Lambda_\theta^C$.
The scaling in \eqref{eq:bounded-total-merge} permits a bounded average rate
per candidate pair.  Merge events cannot cause cardinality explosion because
each one removes a component, although evaluating all logits still costs
$O(N^2)$.

The measure-level reverse kernel is exact, but its coordinate merge rate can
be unbounded, for example near $s=T$ under a hard endpoint-cardinality
restriction or as a parent mass approaches zero when the density ratio does
not cancel the $1/m$ factor.  A finite $k_{\max}$ is then a regularized
implementation and need not contain the exact bridge.  The consistency
theorem applies only under its stated realizability assumption.  If exact hard final-count support is required,
one may add a target-count bridge; we leave that extra conditioning variable
out of the present model.

For large sets, construct a symmetric candidate graph $\mathcal C(X)$ using
$k$ nearest neighbors or learned locality, and normalize
\eqref{eq:pair-softmax} only over $(a,b)\in\mathcal C(X)$.  In this sparse
mode replace \eqref{eq:bounded-total-merge} by
\begin{equation}
  \Lambda_\theta^C(X,s)
  =k_{\max}(s)|\mathcal C(X)|
   \sigma\!\left(g_\theta(g,s,\ell_M,y)\right).
  \label{eq:sparse-total-merge}
\end{equation}
This preserves the bounded average rate per scored candidate and reduces
pair scoring to $O(|\mathcal C|)=O(Nk)$.  It is an approximation unless the
true reverse assigns zero rate to omitted pairs.
During likelihood training, every observed sibling pair must have nonzero
model support.  This can be enforced by using all pairs, by constructing an
equivariant candidate graph with guaranteed support, or by mixing a sparse
proposal with a positive full-pair fallback.  Merely adding the target pair
only during training would define a different model and is not valid.

\begin{theorem}[Permutation invariance of the generated law]
\label{thm:permutation}
Suppose the set encoder is permutation equivariant, the pooled summary is
invariant, the pair logit is symmetric, and the leaf-count and birth-mark
heads depend only on invariant summaries.  In sparse mode, also suppose the
candidate construction and its tie-breaking are permutation equivariant and
retain the stated support.  Starting from $\emptyconfig$, the
distribution of the reverse process is invariant under every relabeling of
component indices at every time.
\end{theorem}

\begin{proof}
The empty state and sampled scalar $L$ are invariant.  Assume the current law
is invariant.  The fixed birth timing rule and birth-mark law are unchanged
by relabeling, so a birth produces the same unordered law.  Merge logits
permute in the same way as their pairs, while the softmax denominator is
invariant; hence pair selection is equivariant.  The merge map itself is
symmetric in its two inputs.  The deterministic flow acts identically on
every component.  Induction over the finite event sequence proves invariance
at every time.
\end{proof}

\begin{proposition}[Non-explosion of the learned sampler]
\label{prop:reverse-nonexplosion}
Assume $k_{\max}$ is bounded on $[S_N,T]$.  For every sampled
$L\leq N_{\max}$, the reverse process has at most $L$ births and
$\max\{L-1,0\}$ merges, hence at most $2N_{\max}-1$ jumps.
\end{proposition}

\begin{proof}
The conditional timing rule produces exactly $L$ births.  Stage C contains
no births, and each merge lowers cardinality by one, so it can contain at
most $\max\{L-1,0\}$ merges.  The bound follows.  Boundedness of $k_{\max}$
ensures every waiting-time calculation before the finite merge count is
well defined.
\end{proof}

\section{Algorithms}
\label{sec:algorithms}

\subsection{Closed-form flow update}

We use the following routine in all algorithms.  For a forward segment from
$t_0$ to $t_1$, compute
\begin{equation}
  G=\int_{t_0}^{t_1}\gamma(u)\dd u,
  \qquad r_1=r_0e^{-G},
  \qquad c_1=M\frac{r_1}{1+r_1},
  \label{eq:flow-update-forward}
\end{equation}
where $r_0=c_0/(M-c_0)$.  Set $m_i\leftarrow m_i(c_1/c_0)$ and
$v\leftarrow M-c_1$.  For a reverse segment from $s_0$ to $s_1$, use
\begin{equation}
  G_R(s_0,s_1)
  =\int_{s_0}^{s_1}\gamma(T-u)\dd u
  =\int_{T-s_1}^{T-s_0}\gamma(t)\dd t.
  \label{eq:reverse-flow-integral}
\end{equation}
Then replace $e^{-G}$ by $e^{+G_R}$.  Both updates preserve total mass up to
floating-point roundoff.  Setting $v=M-\sum_i m_i$ after each update removes
accumulated roundoff.
Before evaluating the ratio, handle the two boundary cases.  If $c_0=0$,
return the unchanged empty state.  If $v_0=0$, return the unchanged state.
In both cases the flow vector field is zero, and these branches avoid $0/0$
or an infinite odds ratio.

\begin{algorithmblock}{Algorithm 1: Exact forward path simulation}
\begin{enumerate}
  \item \textbf{Input:} data state $X_0$, horizon $T$, switch time $\tau$, split rate $\beta$, split density $\kappa$, evaporation schedule $\gamma$.
  \item Set $t\leftarrow0$, $X\leftarrow X_0$, and initialize an empty event
  record.  Give every component a persistent lineage identifier using the
  canonical convention of Section~\ref{sec:learning}.
  \item \textbf{Fragmentation stage:} while $t<\tau$:
  \begin{enumerate}
    \item If $N(X)=0$ or $N(X)=N_{\max}$, apply the forward flow to $\tau$,
    set $t\leftarrow\tau$, and stop this stage.
    \item Draw $E\sim\operatorname{Exp}(1)$ and seek the smallest $t'>t$
    satisfying
    \begin{equation*}
      \int_t^{t'}N(X)\beta(u)\dd u=E.
    \end{equation*}
    If no solution exists before $\tau$, apply the flow to $\tau$, set
    $t\leftarrow\tau$, and stop this stage.  For constant $\beta>0$, this is
    $t'=t+E/[N(X)\beta]$; if $\beta=0$, no solution exists.
    \item Apply the flow from $t$ to $t'$ and set $t\leftarrow t'$.
    \item Sample a parent uniformly and sample $(\rho,u)\sim\kappa_t$.
    \item Apply $X\leftarrow\splitmap_i(X;\rho,u)$.  Record the time, parent identifier, child identifiers, and split mark.
  \end{enumerate}
  \item Set $t\leftarrow\tau$ if the fragmentation loop ended at the stage
  boundary.
  \item \textbf{Deletion stage:} for every component alive at $\tau$, independently draw $D_i\sim\operatorname{Uniform}(\tau,T)$.  Sort the deletion times.
  \item Traverse the sorted times.  Before each deletion, apply the forward
  flow from $t$ to $D_i$, set $t\leftarrow D_i$, record the component state
  $(m_i,z_i)$ immediately before deletion, and apply
  $X\leftarrow\deathmap_i(X)$.
  \item Apply the flow from $t$ to $T$ if needed, set $t\leftarrow T$, and
  set $X_T=\emptyconfig$.
  \item \textbf{Return:} the event record and the reversible flow segments.
\end{enumerate}
\end{algorithmblock}

The deletion times can be sampled independently because Stage E contains no splits and every component has the same hazard \eqref{eq:death-hazard}.  Sorting them is equivalent to simulating competing deletion clocks.

\begin{algorithmblock}{Algorithm 2: One stochastic training step}
\begin{enumerate}
  \item Sample a minibatch of data.  For normalized raw weights, also draw
  $\xi\sim q_\epsilon$ and form the augmented state $X_0$ using
  \eqref{eq:reservoir-dequantization}.  Simulate one forward history for
  each $X_0$ using Algorithm~1.
  \item Reverse each history: reverse the event order, turn deletions into births with their recorded $(m,z)$ marks, and turn splits into merges of the recorded sibling lineages.
  \item Record the leaf budget $L=N(X_\tau)$ and evaluate:
  \begin{itemize}
    \item the leaf-count term $\log\Pi_\theta(L\mid M,y)$;
    \item for each birth, the absolute-mass log density
    \begin{equation*}
      \log A_\theta^{(m)}(m,z\mid X,s,L)
      =\log A_\theta^{(\omega)}(m/v,z\mid X,s,L)-\log v;
    \end{equation*}
    \item for each merge,
    $\log\Lambda_\theta^C(X,s)+\log\pi_\theta(a,b\mid X,s)$.
  \end{itemize}
  In sparse mode, enforce the nonzero-support requirement stated after
  \eqref{eq:sparse-total-merge}; otherwise an observed merge can have
  infinite negative log likelihood.
  \item Estimate the merge survival integral in \eqref{eq:path-loss}.
  Either sample uniform times as in \eqref{eq:uniform-integral-estimator},
  or use the interval-length-weighted stratification in
  \eqref{eq:stratified-integral-estimator}.  Birth times have a fixed law and
  require no learned survival term.
  \item Average \eqref{eq:path-loss} over the minibatch, differentiate it, and update $\theta$.
  \item If a numerical ELBO is wanted, also accumulate the prescribed
  forward density, joint Jacobian, fixed birth-time density, and any
  dequantization terms.  None is needed for the gradient with respect to
  $\theta$.
\end{enumerate}
\end{algorithmblock}

\begin{algorithmblock}{Algorithm 3: Event-driven reverse sampling by thinning}
\begin{enumerate}
  \item \textbf{Input:} trained model, $M$, condition $y$, horizon $T$,
  switch time $\tau$, and valid upper envelopes for merge rates.
  \item Set $s\leftarrow0$ and $X\leftarrow\emptyconfig$.
  Sample $L\sim\Pi_\theta(\cdot\mid M,y)$.
  \item \textbf{Nucleation stage, $0<s<T-\tau$:}
  \begin{enumerate}
    \item If $L=0$, apply reverse growth from $0$ to $T-\tau$ (the empty
    state remains unchanged), set $s\leftarrow T-\tau$, skip the remaining
    nucleation substeps, and continue to the coagulation stage.
    \item Draw $L$ independent times from
    $\operatorname{Uniform}(0,T-\tau)$ and sort them as
    $0<s_1<\cdots<s_L<T-\tau$.
    \item For $j=1,\ldots,L$, apply reverse growth from the current time to
    $s_j$, set $s\leftarrow s_j$, sample $(\omega,z)$ from
    \eqref{eq:birth-mixture}, set $m=\omega v$, and apply
    $X\leftarrow\birthmap_{m,z}(X)$.
    \item After the last birth, apply reverse growth to $T-\tau$ and set
    $s\leftarrow T-\tau$.
  \end{enumerate}
  \item \textbf{Coagulation stage, $T-\tau<s<T$:}
  \begin{enumerate}
    \item If $N(X)<2$, apply reverse growth to $T$ and stop.  If a sparse
    candidate set is empty, fall back to the full pair set by default, or
    advance only to the next scheduled candidate-refresh time.  Stop early
    only if emptiness is known to remain valid until $T$.
    \item Choose an interval $I=[s,s_I]$ and a valid upper envelope
    $\overline K$ for $\Lambda_\theta^C$ along its deterministic flow.  For
    full pair scoring, a direct choice is
    \begin{equation*}
      \overline K=\binom{N(X)}{2}\sup_{u\in I}k_{\max}(u).
    \end{equation*}
    In sparse mode, replace the binomial factor by a valid upper bound on
    $|\mathcal C(X_u)|$ over $u\in I$.
    If $\overline K=0$, flow to $s_I$, set $s\leftarrow s_I$, and recompute;
    return to the coagulation-stage test, and stop only when $s_I=T$.
    \item Draw $E_K\sim\operatorname{Exp}(\overline K)$.  If
    $s+E_K\geq s_I$,
    flow to $s_I$, set $s\leftarrow s_I$, and return to the coagulation-stage
    test.  Otherwise flow to $s+E_K$ and set $s\leftarrow s+E_K$.
    \item Accept with probability
    $\Lambda_\theta^C(X,s)/\overline K$.  On acceptance, sample an unordered
    pair from $\pi_\theta(a,b\mid X,s)$ and apply
    $X\leftarrow\mergemap_{ab}(X)$.  Whether accepted or rejected, recompute
    or revalidate the next envelope and repeat the coagulation-stage test.
  \end{enumerate}
  \item Return the unordered component set and report the remaining
  reservoir $v$.  For a nonempty normalized-weight output, decode
  $w_i=m_i/\sum_jm_j$; treat the empty set as a separate outcome.
\end{enumerate}
\end{algorithmblock}

The thinning algorithm assumes a finite envelope on each interval.  If an
unbounded bridge schedule is deliberately used near $T$, simulate it through
its integrated hazard or a time change; a single finite envelope extending
to $T$ does not exist.

When analytic envelopes are inconvenient, a time grid may replace merge
thinning.  At the left endpoint $(X,s)$, evaluate all merge rates and set
$\Lambda=K_\theta^{\mathrm{tot}}(X,s)$.  Fix a step-size tolerance
$\rho_{\mathrm{grid}}\in(0,1)$ and choose $\Delta\leq T-s$.  When
$\Lambda>0$, also require
$\Lambda\Delta\leq\rho_{\mathrm{grid}}$ and use
\begin{equation}
  p_{\mathrm{event}}=1-e^{-\Lambda\Delta}.
  \label{eq:grid-event-probability}
\end{equation}
Draw a Bernoulli variable with this probability.  If it is one, draw and
retain a pair from $K_\theta(a,b\mid X,s)/\Lambda$.  Apply exact reverse
growth from $s$ to $s+\Delta$, set $s\leftarrow s+\Delta$, and then merge the
retained pair if an event was drawn.  Thus rates and pair probabilities are
frozen at the left endpoint and at most one event occurs per step.  This is a
first-order approximation and converges as $\max\Delta\to0$.  If a
multiple-event tau-leap is used instead, events must be applied sequentially
because each merge changes both the available pairs and their rates.  The
grid changes timing accuracy but not mass conservation.  When $\Lambda=0$,
set $p_{\mathrm{event}}=0$, apply only the growth update over $\Delta$, and
do not evaluate $K_\theta/\Lambda$.

\subsection{Computational complexity}

Let $C_{\mathrm{birth}}(N)$ be the cost of one birth-mark evaluation at
cardinality $N$, and let $C_{\mathrm{enc}}(N)$ be the cost of one set-encoder
evaluation.  Let $N_\ell$ be the active cardinality at birth evaluation
$\ell$.  Let $P_C$ be the number of merge-network evaluations, including
accepted events, rejected thinning proposals, and survival-integral queries,
and let $N_j$ be the active cardinality at its $j$th evaluation.  Sorting
birth times costs $O(L\log L)$.  With full pair scoring, the total sampling
and rate-evaluation cost is
\begin{equation}
  O\!\left(
    L\log L
    +\sum_{\ell=1}^{L}C_{\mathrm{birth}}(N_\ell)
    +\sum_{j=1}^{P_C}
      \left[C_{\mathrm{enc}}(N_j)+N_j^2\right]
  \right).
  \label{eq:full-complexity}
\end{equation}
A sparse $k$-neighbor candidate graph replaces $N_j^2$ by
$|\mathcal C(X_j)|=O(N_jk)$.  Batched survival evaluations dominate training
when many quadrature times are used; caching the fixed reversed path states
and batching all queried times reduces overhead.

\section{Further theoretical properties}
\label{sec:properties}

\subsection{Endpoint error is controlled by path error}

Let $\mathbb P^*$ be the exact augmented reverse path law obtained from
Theorem~\ref{thm:reverse-kernel}, and let $\mathbb P_\theta$ be the learned
reverse path law.  Both begin at $\emptyconfig$ and first sample their
respective leaf budgets.  Write $P_{\mathrm{data}}$ and $P_\theta^T$ for
their endpoint probability measures.  These are induced by the augmented
data law and learned endpoint densities, respectively.

\begin{theorem}[Path-to-endpoint error bound]
\label{thm:endpoint-bound}
If $\mathbb P^*\ll\mathbb P_\theta$, then
\begin{equation}
  \KL(P_{\mathrm{data}}\,\|\,P_\theta^T)
  \leq
  \KL(\mathbb P^*\,\|\,\mathbb P_\theta).
  \label{eq:data-processing-kl}
\end{equation}
Consequently,
\begin{equation}
  \TV(P_{\mathrm{data}},P_\theta^T)
  \leq
  \sqrt{\frac12\KL(\mathbb P^*\,\|\,\mathbb P_\theta)}.
  \label{eq:pinsker-endpoint}
\end{equation}
\end{theorem}

\begin{proof}
The endpoint is a measurable function of the full path.  Relative entropy cannot increase under a measurable map, which proves \eqref{eq:data-processing-kl}.  Equation~\eqref{eq:pinsker-endpoint} is Pinsker's inequality.
\end{proof}

For normalized observations, applying the decoder
\eqref{eq:normalized-decoder} is another measurable map.  The same
data-processing argument therefore also bounds the mismatch between the
decoded nonempty endpoint laws, with the empty outcome retained as its
separate atom.

When the exact and learned laws use the same deterministic flow, the same
conditional birth-time law, a common record reference measure and support,
and integrable intensities, the path-space KL decomposes into the leaf-count
KL, conditional birth-mark KL terms, and the merge-intensity divergence in
\eqref{eq:intensity-divergence}.  Equivalently, the \emph{population excess}
of the full reverse-record negative log likelihood over its value at the
exact reverse law equals this path KL.  It is this excess risk, not the raw
nonnegative loss value, that upper-bounds the endpoint KL through
Theorem~\ref{thm:endpoint-bound}.  If the flow, timing law, or control law is
also learned, its corresponding path-space term must be added; incompatible
supports give infinite KL.

\subsection{Hard conservation is stronger than a penalty}

\begin{corollary}[Zero structural conservation error]
In exact arithmetic, every sample returned by Algorithm~3 satisfies
\begin{equation}
  v+\sum_i m_i=M.
\end{equation}
This statement holds for every parameter value $\theta$, not only at an optimum.
\end{corollary}

\begin{proof}
Algorithm~3 is a composition of the reverse growth flow, birth maps, and merge maps.  Proposition~\ref{prop:event-conservation} shows that each map is conservative.
\end{proof}

This is why the proposed method does not include a mass-conservation loss.  A penalty can make violations small on average; it cannot provide the above structural guarantee.

\subsection{Conditioning on varying total mass}

If total mass varies across observations, first model
\begin{equation}
  M\sim p_\vartheta(M\mid y)
\end{equation}
with a one-dimensional density, or treat $M$ as observed conditioning information.  Conditional on the sampled or supplied $M$, the entire coagulation--fragmentation process remains conservative.  Mixing examples with different $M$ without either step would make the claimed conservation law ill-defined.
The parameterization above supplies $\ell_M=\log(M/M_{\mathrm{ref}})$ to the encoder, birth, pair,
and total-rate heads.  Omitting that input would restrict the model to
scale-invariant conditional laws even if $M$ were sampled separately.
In the first case, training adds
\begin{equation}
  \mathcal L_M=-\log p_\vartheta(M\mid y),
\end{equation}
and Algorithm~3 samples $M\sim p_\vartheta(\cdot\mid y)$ before initializing
$\emptyconfig$.  In the experiments proposed here, we instead condition on
or normalize $M$ so that this extra head is not required.

\subsection{Feature diffusion is outside the present model}

All theorems, likelihoods, and algorithms in this paper assume deterministic
feature motion between jumps.  Adding Brownian feature noise would require a
reverse-score or Girsanov term and an SDE solver.  That is a possible future
extension, but it is not part of the model analyzed here.

\section{Extension: locally controlled evaporation}
\label{sec:local-control}

\subsection{What the hair-dryer observation adds}

The mirror episode suggests more than a global evaporation schedule.  A
hair dryer clears a small region, the user sees which regions remain covered,
and the dryer is moved again.  The useful abstract properties are therefore
\emph{locality}, \emph{sequential action}, and \emph{feedback}.  This section
adds those properties without changing the split, merge, deletion, or birth
maps of the core model.

There is also an important distinction.  On a real mirror, the reflected
face already exists behind the droplets.  Literal mirror clearing is thus a
restoration or de-occlusion problem, not unconditional generation.  We use
the hair dryer as motivation for a spatial control mechanism.  We do not
claim that the hidden face and the generated weighted set are the same
mathematical object.

\subsection{A spatial control field}

Some coordinates of $z_i$ need not be spatial.  Let
$x_i^{\mathrm{sp}}=\operatorname{pos}(z_i)\in\R^{q_s}$ be a specified spatial projection,
and let $a_t:\R^{q_s}\to[0,1]$ describe where the dryer acts at forward
time $t$.
A convenient finite representation uses $K_{\mathrm{basis}}$ nonnegative
spatial basis functions satisfying
$\sum_{k=1}^{K_{\mathrm{basis}}}\psi_k(x)\leq1$:
\begin{equation}
  a_t(x)=\sum_{k=1}^{K_{\mathrm{basis}}}a_{t,k}\psi_k(x),
  \qquad a_{t,k}\in[0,1].
  \label{eq:local-control-basis}
\end{equation}
The functions may be image patches, radial kernels, or cells of a spatial
grid.  The coefficient vector determines the location and strength of the
current action.

For component $i$, let
\begin{equation}
  \lambda_i(t,X,a_t)
  =\lambda_0(t)+\lambda_1(t)a_t(x_i^{\mathrm{sp}})\geq0,
  \label{eq:local-evaporation-rate}
\end{equation}
where $\lambda_0(t),\lambda_1(t)\geq0$, or use any bounded nonnegative
permutation-equivariant rate.  Replace the
global flow by
\begin{align}
  \frac{\dd m_i}{\dd t}
  &=-\lambda_i(t,X_t,a_t)\frac{v}{M}m_i,
  \label{eq:controlled-flow-mi}\\
  \frac{\dd v}{\dd t}
  &=\frac{v}{M}\sum_{i=1}^{N(X_t)}
    \lambda_i(t,X_t,a_t)m_i,
  \label{eq:controlled-flow-v}\\
  \frac{\dd z_i}{\dd t}&=0.
\end{align}
The original model is recovered by setting
$\lambda_i(t,X,a)=\gamma(t)$ for every component.  Unlike the global flow,
the controlled flow generally changes relative component masses and does
not reduce to one scalar odds equation.

\begin{proposition}[Conservation and positivity under local control]
\label{prop:controlled-conservation}
Assume the action path is piecewise continuous.  Assume the rates
$\lambda_i$ are bounded, nonnegative, and uniformly Lipschitz in the state
on the compact mass simplex, uniformly over time and admissible actions.
Then the controlled ODE has a unique solution on every finite time interval.  If initially
$m_i\geq0$, $v\geq0$, and $v+\sum_i m_i=M$, these properties hold for all
later times.  If the rates are permutation equivariant, so is the flow.
\end{proposition}

\begin{proof}
On each interval where the action is continuous, the regularity assumptions
give a unique local solution; concatenating the intervals handles action
jumps.  The state remains
in the compact mass simplex, so the bounded vector field cannot explode and
the solution extends through every finite interval.  At $m_i=0$, the right
side of \eqref{eq:controlled-flow-mi} is zero.  At $v=0$, all mass
derivatives in \eqref{eq:controlled-flow-mi}--\eqref{eq:controlled-flow-v}
are zero.  Thus no coordinate can cross below zero.  Finally,
\begin{equation*}
  \frac{\dd}{\dd t}\left(v+\sum_i m_i\right)
  =\frac{v}{M}\sum_i\lambda_i m_i
   -\frac{v}{M}\sum_i\lambda_i m_i=0.
\end{equation*}
The total is therefore $M$.  Relabeling the components only relabels an
equivariant vector field, which proves permutation equivariance.
\end{proof}

\subsection{Conditional time reversal}

The cleanest exact construction treats the complete action path
$a_{0:T}$ as known conditioning information.  Let $P_t^a$ and $Q_t^a$
denote the forward marginal and jump kernel conditional on this path.  The
reverse jumps still satisfy the same flux identity,
\begin{equation}
  P_t^a(\dd Y)\overline Q_{T-t}^a(Y,\dd X)
  =P_t^a(\dd X)Q_t^a(X,\dd Y).
  \label{eq:controlled-flux-reversal}
\end{equation}
Between jumps, the reverse traverses the controlled ODE backward:
\begin{align}
  \frac{\dd m_i}{\dd s}
  &=+\lambda_i(T-s,Y_s,a_{T-s})\frac{v}{M}m_i,
  \label{eq:controlled-reverse-mi}\\
  \frac{\dd v}{\dd s}
  &=-\frac{v}{M}\sum_i
    \lambda_i(T-s,Y_s,a_{T-s})m_i.
  \label{eq:controlled-reverse-v}
\end{align}
Thus the reverse process performs stronger growth in the spatial regions
that were evaporated more strongly in the forward direction.

The closed-form flow Jacobian \eqref{eq:flow-jacobian} must be replaced by
an ODE Jacobian.  Let $b^a(m,t)$ be the vector formed by the right sides of
\eqref{eq:controlled-flow-mi}, with $v=M-\sum_i m_i$.  For a no-jump
segment from $r$ to $t$, additionally assume that $b^a(\cdot,t)$ is $C^1$
on the mass-simplex interior and has integrable divergence along the
trajectory.  Liouville's formula then gives
\begin{equation}
  \log\left|\det D\Phi_{r,t}^{a}(m_r)\right|
  =\int_r^t \nabla_m\!\cdot b^a(m_u,u)\dd u.
  \label{eq:controlled-flow-jacobian}
\end{equation}
Equivalently, the multiplicative flow factor used in a change of variables is
\begin{equation}
  J_{\mathrm{flow}}^a(r,t)
  =\exp\!\left\{
    \int_r^t \nabla_m\!\cdot b^a(m_u,u)\dd u
  \right\}.
  \label{eq:controlled-flow-factor}
\end{equation}
The log factor can be integrated together with the state.  When each
$\lambda_i$ depends on $(t,x_i^{\mathrm{sp}},a_t)$ but not on the masses,
\begin{equation}
  \nabla_m\!\cdot b^a(m,t)
  =\frac{1}{M}\sum_i\lambda_i(t,x_i^{\mathrm{sp}},a_t)(m_i-v).
  \label{eq:controlled-flow-divergence}
\end{equation}
For a mass-dependent neural rate, automatic differentiation gives the full
divergence.  Equation~\eqref{eq:controlled-flow-divergence} treats the action
on each segment as fixed conditioning information.  If the action is
recomputed differentiably from the current state, the divergence must also
differentiate through that controller.

The reverse merge-rate factor $1/m$ in \eqref{eq:exact-merge-rate} is
unchanged.  Separately, every forward split still contributes its parent
mass $m$ to the joint-history Jacobian in
\eqref{eq:joint-history-jacobian}.  Each global-flow factor in that Jacobian
is replaced by $J_{\mathrm{flow}}^a$ from
\eqref{eq:controlled-flow-factor}.  For one prescribed control path, the
algebraic form of the trainable jump loss \eqref{eq:path-loss} is unchanged,
but $A_\theta$ and $K_\theta$ must condition on that path, or on a sufficient
current control-state encoding.  If the control or controlled-flow
parameters are learned jointly with the generator, the action density and
controlled-flow Jacobian are parameter-dependent and must also be included
in the objective.

\subsection{Feedback and adaptive stopping}

An open-loop controller fixes $a_{0:T}$ before the process starts.  The
mirror story instead suggests a closed-loop controller.  At decision time
$t_k$, let $o_k$ be an observation or uncertainty summary and sample
\begin{equation}
  a_k\sim\pi_\varphi(\,\cdot\mid X_{t_k},o_k,t_k,h_k^{\mathrm{ctrl}}),
  \label{eq:control-policy}
\end{equation}
where $h_k^{\mathrm{ctrl}}$ is optional controller memory.  We take $o_k$
and the memory update to be deterministic functions of the augmented state
and past actions.  If either is stochastic, its transition density must also
be included in the history likelihood.  The sampled coefficient vector is
held fixed on $[t_k,t_{k+1})$, or interpolated by a stated deterministic
rule, to define the continuous field $a_t$ used by the ODE.

The process is Markov after augmenting its state with
$(o_k,h_k^{\mathrm{ctrl}},a_k)$.  There is one extra point for fixed decision
times: these are predictable transitions and are not covered by the
no-predictable-jump condition in
Assumption~\ref{ass:reversal-regularity}.  Let $Z_k^-$ and $Z_k^+$ be the
augmented states immediately before and after decision $t_k$, and let
$R_k(Z_k^-,\dd Z_k^+)$ be the forward scheduled-transition kernel.  Its
pre- and post-decision laws are denoted by
$\mathsf P_{t_k^-}$ and $\mathsf P_{t_k}$.  Its exact reverse is the
discrete Bayes kernel $\overline R_k$ defined by
\begin{equation}
  \mathsf P_{t_k}(\dd Z^+)\overline R_k(Z^+,\dd Z^-)
  =\mathsf P_{t_k^-}(\dd Z^-)R_k(Z^-,\dd Z^+).
  \label{eq:scheduled-control-reversal}
\end{equation}
Between decision times, the continuous flow and ordinary jumps use
\eqref{eq:controlled-flux-reversal}.  Hazard-driven random decision times
can instead be included in the ordinary augmented jump kernel.
Noninvertible memory updates cause no problem for
\eqref{eq:scheduled-control-reversal}, but the forward history must retain
the predecessor state and any sampled noise used to train its reverse.

Let $J$ denote the number of control decisions in a realized trajectory.
Three probability conventions must be kept separate:
\begin{enumerate}
  \item A prescribed action path is conditioning information and contributes
  no action-density term.
  \item If $\pi_\varphi$ is the forward controller used in the variational
  corruption law, then
  \begin{equation}
    \log q_\varphi(a_{1:J},H\mid X)
    \supset
    \sum_k\log\pi_\varphi
    (a_k\mid X_{t_k},o_k,t_k,h_k^{\mathrm{ctrl}}),
    \label{eq:control-path-density}
  \end{equation}
  and this term enters the ELBO with a minus sign through $-\log q_\varphi$.
  \item If actions are latent variables of the generative model, parameterize
  a reverse scheduled kernel $\overline R_{\omega,k}$ satisfying
  \begin{equation*}
    Z_k^-\sim\overline R_{\omega,k}(Z_k^+,\dd Z_k^-),
  \end{equation*}
  or use a hazard-driven reverse action kernel.  Its eventwise log density
  $\sum_k\log\overline r_{\omega,k}(Z_k^-\mid Z_k^+)$ enters the generative
  path likelihood with a plus sign.
\end{enumerate}
An online forward controller cannot be used during generation merely by
``reversing its actions,'' because no realized forward action path exists at
sampling time.  Its action transitions must themselves be reversed through
the augmented-state flux identity or modeled by the separate reverse kernel
in item 3.

Let $U_\psi(o_k)\geq0$ measure how much of the target remains unresolved.
A simple stopping time is
\begin{equation}
  \tau_{\mathrm{stop}}
  =\min\!\left\{T,\inf\{t_k:U_\psi(o_k)\leq\epsilon\}\right\},
  \label{eq:control-stopping-time}
\end{equation}
where $\inf\varnothing=+\infty$.
The controller can be learned from action demonstrations using the negative
of \eqref{eq:control-path-density}.  It may instead be optimized with the
cost
\begin{equation}
  \mathcal J_{\mathrm{ctrl}}(\varphi)
  =\E\!\left[
    U_\psi(o_J)
    +\lambda_a\sum_{k=1}^{J}\operatorname{Cost}(a_k)
    +\lambda_{\mathrm{step}}J
  \right].
  \label{eq:control-objective}
\end{equation}
The first term rewards revealing the target, the second penalizes dryer
energy or spatial coverage, and the third penalizes long action sequences.
This feedback objective is useful for conditional restoration or inpainting.
For unconditional generation, the fixed horizon $T$ remains the simpler and
better-defined stopping rule.  Equation~\eqref{eq:control-stopping-time}
therefore defines a separate control or restoration objective; it is not
silently part of the fixed-horizon generative model.  If adaptive stopping is
used inside that model, the stopping decision needs an explicit probability
or hazard and a random terminal law.  Both must appear in the forward
variational and reverse generative path densities.

\subsection{Implementation}

\begin{algorithmblock}{Algorithm 4: Controlled-flow modification}
\begin{enumerate}
  \item Choose spatial basis functions
  $\psi_{1:K_{\mathrm{basis}}}$.  Select one action convention: a prescribed
  path; a forward variational controller together with a learned reverse
  action kernel; or a generative action prior with its reverse controller.
  \item In Algorithm~1, keep the same split and deletion events.  Between
  events, integrate \eqref{eq:controlled-flow-mi}--\eqref{eq:controlled-flow-v}
  and, when a numerical ELBO is required, simultaneously integrate the
  divergence in \eqref{eq:controlled-flow-jacobian}.  Record every action,
  observation, and controller-memory state used by the forward law.
  \item For a prescribed path, reverse its time ordering and use $a_{T-s}$.
  For latent or online actions at fixed decision times, train and sample
  $\overline R_{\omega,k}$ using
  \eqref{eq:scheduled-control-reversal}, and add
  $\sum_k\log\overline r_{\omega,k}$ to the reverse-record log likelihood.
  For hazard-driven decisions, use the augmented continuous-time reverse
  kernel.  Do not replay a forward action path that does not exist at
  generation time.
  \item In Algorithm~3, replace each closed-form growth update by numerical
  integration of
  \eqref{eq:controlled-reverse-mi}--\eqref{eq:controlled-reverse-v}.
  Birth and merge sampling are otherwise unchanged.
  \item A prescribed path adds no action-density term.  A forward variational
  controller contributes through $-\log q_\varphi$, whereas a generative
  action prior or reverse policy contributes through $+\log p$.  Use
  \eqref{eq:control-stopping-time} only for a modeled conditional revealing
  task; otherwise retain the fixed horizon $T$.
\end{enumerate}
\end{algorithmblock}

Use a positivity-preserving, mass-conserving ODE integrator with step
rejection.  Recomputing $v=M-\sum_i m_i$ solely to remove floating-point
roundoff is harmless.  Material clipping or projection, however, changes the
numerical map and invalidates the continuous Liouville Jacobian in
\eqref{eq:controlled-flow-jacobian}.  If such a discrete solver is used, the
implementation must instead compute the Jacobian of its actual discrete map.
The local model is slower than the scalar closed-form flow, so the global
model should remain the default unless spatial control yields a measurable
benefit.

\subsection{Airflow, departure, and the modeled boundary}

A physical hair dryer can also translate droplets or push them off the
mirror.  Deterministic translation of the spatial coordinates may be
represented by
$\dd x_i^{\mathrm{sp}}/\dd t
=u_{\mathrm{air}}(t,x_i^{\mathrm{sp}},a_t)$.  If water leaves the modeled
region, transferring it directly to $v$ would incorrectly call it local
vapor.  Introduce an escaped reservoir $e$ instead and conserve
\begin{equation}
  v+e+\sum_i m_i=M.
  \label{eq:extended-mass-balance}
\end{equation}
A departure event removes component $i$ and adds $m_i$ to $e$; its reverse
is an insertion drawing mass from $e$.  These transport and departure events
are not needed for the minimal local-evaporation extension and are left out
of the main experiments.

\subsection{Experiments specific to this extension}

The controlled model should first be tested on synthetic spatially
heterogeneous sets or a conditional de-occlusion task.  Compare four matched
controllers: global uniform action, a fixed raster scan, an uncertainty-based
oracle, and a learned feedback policy.  Report generation or restoration
quality together with the action energy
$\int_0^T\|a_t\|_1\dd t$, number of control decisions, wall-clock time, and
mass residual.  The extension is useful only if it improves quality or
reduces computation at a matched action budget.

\section{Experimental plan (results pending)}
\label{sec:experiments}

\begin{center}
\fbox{\parbox{0.92\linewidth}{
\textbf{Status.} No experiment in this section has been run.  The hypotheses,
baselines, and metrics below are a prospective evaluation plan.  Numerical
tables are intentionally blank, and the text makes no empirical claim.
}}
\end{center}

\subsection{Research questions}

\begin{description}[leftmargin=1.0cm,style=nextline]
  \item[H1: Reverse correctness.] On a finite process with exactly computable marginals, learned reverse rates approach the analytical rates.
  \item[H2: Benefit of coagulation.] At matched architecture and compute,
  pairwise merging captures pair and cluster statistics better than a
  single-component insertion model with no pair-merge operation.
  \item[Check 1: Conservation.] Verify as an implementation check that the
  sampler preserves total mass to numerical precision; this is a theorem,
  not an empirical hypothesis.
  \item[H3: Cardinality.] After sampling an auxiliary leaf budget and then
  merging leaves, the model matches the held-out \emph{final} cardinality
  distribution.  This does not claim that the method has no count model.
  \item[H4: Efficiency.] Sample quality as a function of wall-clock time, event count, and peak active cardinality is competitive with trans-dimensional baselines.
  \item[H5: Local control.] On a spatially heterogeneous conditional task,
  the extension in Section~\ref{sec:local-control} improves quality or
  reduces action cost relative to uniform evaporation and a fixed scan.
\end{description}

\subsection{Benchmarks}

\begin{enumerate}
  \item \textbf{Finite enumerated process.}  Use integer masses, a small
  feature lattice, and small $N_{\max}$.  Compute $P_t$ and the exact reverse
  generator by solving the finite Kolmogorov equation.  This tests probability-flux reversal
  and unordered pair counting, but not the continuous $1/m$ Jacobian.
  \item \textbf{Continuous synthetic weighted sets.}  Sample a known
  hierarchical fragmentation process with controllable split imbalance,
  spatial interaction strength, and cardinality.  Retain the ground-truth
  genealogy only for evaluation.  Separately test the continuous split
  change of variables by Monte Carlo integration to verify the $1/m$ factor.
  \item \textbf{One real scalar-weighted set domain.}  Use a dataset in
  which the additive scalar in the present theory is measured rather than
  invented.  Two compatible candidates are calorimeter-hit point clouds
  with deposited energy as $m_i$, conditioned on total deposited energy,
  and instance-segmented droplet microscopy with calibrated droplet volume
  as $m_i$.
  Vector four-momentum conservation is outside the present scalar theory and
  is left to future work.
\end{enumerate}

\subsection{Baselines}

\begin{itemize}
  \item Jump Diffusion Models with the same set encoder \cite{campbell2023transdimensional}.
  \item Existence-Field Diffusion on spatial point-process tasks \cite{pan2026existence}.
  \item A fixed- or cardinality-conditioned diffusion baseline.
  \item An autoregressive split-tree model when a branching domain is used.
  \item The analytical reverse generator on the finite benchmark.
  \item \textbf{Critical in-house ablation:} a birth-only
  single-component insertion process using the same backbone, parameter
  count, training paths, and compute budget, but no pairwise merging.
\end{itemize}

\subsection{Metrics}

Report:
\begin{itemize}
  \item reverse-rate relative error and event calibration on the exact task;
  \item birth-mark negative log likelihood and merge-pair accuracy;
  \item total variation or Hellinger distance between cardinality distributions;
  \item Wasserstein distance for component-mass marginals;
  \item pair-correlation or Ripley-$K$ error for spatial structure;
  \item MMD, energy distance, and a classifier two-sample test on invariant set summaries;
  \item the relative conservation residual
  \begin{equation}
    \epsilon_M=
    \frac{|v+\sum_i m_i-M|}{M+10^{-12}};
    \label{eq:conservation-metric}
  \end{equation}
  \item held-out ELBO on the finite task and on continuous tasks only after
  the joint Jacobian and dequantization terms have been implemented and
  checked; exact likelihood only where it is genuinely computable;
  \item for local control, integrated action energy, spatial coverage,
  control-decision count, and terminal unresolved uncertainty;
  \item wall-clock sampling time, neural evaluations, jump count, average and maximum active cardinality, and peak memory.
\end{itemize}

\subsection{Ablations}

The minimum ablation set is:
\begin{enumerate}
  \item pairwise coagulation versus the birth-only in-house ablation;
  \item exact conservation versus an unconstrained penalty-based model;
  \item learned pair kernel versus uniform and distance-only kernels;
  \item reservoir versus no explicit reservoir;
  \item coupled versus factorized birth marks;
  \item full pair enumeration versus sparse candidate pairs;
  \item several fragmentation strengths and switch times $\tau$;
  \item for the optional extension, global control versus a fixed local scan,
  an uncertainty oracle, and a learned feedback policy.
\end{enumerate}

All models should use identical data splits and comparable encoder capacity.
Report five independent random seeds with 95\% bootstrap confidence intervals
and include quality--compute curves rather than only one operating point.

\subsection{Planned reporting format}

\begin{table}[h]
\centering
\caption{Finite reverse-process validation.  Results pending.}
\begin{tabular}{lccc}
\toprule
Method & Rate error & Mark NLL & Cardinality TV \\
\midrule
Analytical reverse & --- & --- & --- \\
Birth-only in-house ablation & --- & --- & --- \\
Proposed coagulation model & --- & --- & --- \\
\bottomrule
\end{tabular}
\end{table}

\begin{table}[h]
\centering
\caption{Continuous weighted-set generation.  Results pending.}
\begin{tabular}{lcccc}
\toprule
Method & Set distance & Pair-stat. error & $\epsilon_M$ & Time/sample \\
\midrule
Cardinality-conditioned model & --- & --- & --- & --- \\
Jump diffusion & --- & --- & --- & --- \\
Birth-only in-house ablation & --- & --- & --- & --- \\
Proposed coagulation model & --- & --- & --- & --- \\
\bottomrule
\end{tabular}
\end{table}

\section{Limitations}

The method has several clear limitations.  Unrestricted pair scoring is
quadratic in active cardinality.  Sparse candidates reduce cost but may
remove important long-range merges.  The finite cap $N_{\max}$ and the
auxiliary leaf-count model are practical structural choices; a poor cap or
count model can limit support.  Simulated forward histories may also have
high variance, so the fragmentation schedule is a modeling decision rather
than a harmless implementation detail.

The scalar conservation law is useful only when $m_i$ has a defensible
additive meaning.  Normalizing arbitrary scalar scores to sum to one would make
the constraint cosmetic.  The present split map preserves a Euclidean
weighted centroid.  Categorical attributes, molecules with typed atoms, and
systems with vector-valued conservation laws require new event maps and mark
models.  Normalized data also require the reservoir dequantization and
decoder described in Section~\ref{sec:state}.

The prescribed-path local-control extension has an exact conditional
reversal.  A feedback-controlled application additionally needs a
task-specific observation model and the scheduled reverse kernels in
\eqref{eq:scheduled-control-reversal}; Algorithm~4 specifies where these
terms enter but does not choose one universal controller architecture.

The exact terminal construction uses a dedicated deletion stage.  This is
mathematically clean but less physically literal than simultaneous
nucleation, growth, and coagulation.  The condensation language is therefore
an intuition, not a claim that the algorithm implements Hertz--Knudsen
kinetics or an equilibrium free energy.  The general jump-flux identity and
variable-dimensional generation are also prior ideas; the claimed technical
contribution is their conservative unordered split--merge specialization and
the resulting learning and sampling scheme.

Finally, a powerful birth-mark model could explain most of the data while
merges add little.  The birth-only in-house ablation is therefore necessary
to establish that coagulation provides more than an attractive
interpretation.

\section{Conclusion}

We developed a complete theoretical and algorithmic specification of a mass-conserving generative process on variable-size weighted sets.  A prescribed fragmentation--evaporation process reaches an exact all-reservoir state.  Its finite-horizon reverse performs nucleation, conservative growth, and pairwise coagulation.  The reverse rates follow from probability flux, while a path-space ELBO turns simulated corruption histories into a tractable marked point-process loss.  Conservation and permutation invariance hold by construction.

The remaining question is empirical rather than definitional: does pairwise
coagulation improve modeling or efficiency compared with the birth-only
ablation and existing variable-cardinality methods?  Section~\ref{sec:experiments}
specifies the experiments needed to answer that question without
presupposing the result.

\appendix

\section{Derivation of the reservoir dequantization formulas}
\label{app:dequantization-derivation}

This section derives both the Jacobian in
\eqref{eq:dequantization-jacobian} and the likelihood bound in
\eqref{eq:dequantization-bound}.  Write
\begin{equation}
  w_n=1-\sum_{i=1}^{n-1}w_i,
  \qquad
  m_i=M(1-\xi)w_i,
  \quad i=1,\ldots,n.
  \label{eq:appendix-dequantization-map}
\end{equation}
Thus $w_1,\ldots,w_{n-1}$ are the independent simplex coordinates, while
$\xi$ determines the total amount of mass placed in the reservoir.

If $\xi$ were fixed at $\xi_0$, the image would satisfy the extra equality
$\sum_i m_i=M(1-\xi_0)$.  It would therefore have only $n-1$ mass
coordinates inside an $n$-dimensional mass domain, and hence zero
$n$-dimensional Lebesgue measure.  Allowing $\xi$ to vary supplies the
missing direction.  Conversely, for any nonempty interior state,
$\xi=v/M$ and $w_i=m_i/\sum_jm_j$ recover its unique normalized coordinates,
as proved in Proposition~\ref{prop:encode-decode}.

\subsection{Jacobian determinant}

For $n=1$, we have $w_1=1$ and
$\partial m_1/\partial\xi=-M$, so the claimed absolute determinant is
$M=M^1(1-\xi)^0$.  Assume $n\geq2$ for the block calculation below.
Define $a=M(1-\xi)$ and
$\widetilde w=(w_1,\ldots,w_{n-1})^\top$.  For $i<n$,
\begin{equation}
  \frac{\partial m_i}{\partial w_j}=a\delta_{ij},
  \qquad
  \frac{\partial m_i}{\partial\xi}=-Mw_i.
  \label{eq:appendix-jacobian-first-rows}
\end{equation}
Because $m_n=a(1-\sum_{i<n}w_i)=aw_n$, the final row satisfies
\begin{equation}
  \frac{\partial m_n}{\partial w_j}=-a,
  \qquad
  \frac{\partial m_n}{\partial\xi}=-Mw_n.
  \label{eq:appendix-jacobian-last-row}
\end{equation}
Therefore the full Jacobian matrix is
\begin{equation}
  \frac{\partial(m_1,\ldots,m_n)}
       {\partial(w_1,\ldots,w_{n-1},\xi)}
  =
  \begin{pmatrix}
    aI_{n-1} & -M\widetilde w\\
    -a\mathbf 1^\top & -Mw_n
  \end{pmatrix}.
  \label{eq:appendix-dequantization-jacobian-matrix}
\end{equation}
Using the block-determinant identity
$\det\bigl(\begin{smallmatrix}A&B\\C&D\end{smallmatrix}\bigr)
=\det(A)\det(D-CA^{-1}B)$ gives
\begin{align}
  \det J
  &=a^{n-1}
    \left[-Mw_n
    -(-a\mathbf 1^\top)(aI_{n-1})^{-1}(-M\widetilde w)
    \right]
    \nonumber\\
  &=a^{n-1}
    \left[-Mw_n-M\sum_{i=1}^{n-1}w_i\right]
    \nonumber\\
  &=-Ma^{n-1}
   =-M^n(1-\xi)^{n-1}.
  \label{eq:appendix-dequantization-signed-determinant}
\end{align}
Taking the absolute value proves
\begin{equation}
  |\det J|=M^n(1-\xi)^{n-1}.
\end{equation}
The exponent $n-1$ appears because the normalized weight vector has
$n-1$ independent simplex directions, each scaled by $M(1-\xi)$.  The
remaining radial direction is controlled by $\xi$ and contributes the final
factor $M$.  The feature coordinates are unchanged and therefore contribute
a Jacobian factor of one.

\subsection{Raw normalized-data likelihood}

For compact notation, define
\begin{equation}
  f_\theta(\xi;w,z)
  :=p_{\theta,n}(X(\xi))M^n(1-\xi)^{n-1},
  \qquad 0<\xi<1.
  \label{eq:appendix-f-xi}
\end{equation}
Here $p_{\theta,n}$ is the \emph{joint} endpoint density of all masses and
features relative to $\mu_n$; it is not a separately specified density for
$z$.  Under the change of variables above, the model density in coordinates
$(w_{1:n-1},z_{1:n},\xi)$ is
\begin{equation}
  \widetilde p_\theta(w_{1:n-1},z_{1:n},\xi)
  =f_\theta(\xi;w,z).
  \label{eq:appendix-augmented-density}
\end{equation}
The same unordered counting convention is used before and after this change
of variables, so it introduces no additional factorial.  Since $\xi$ is not
part of the raw observation, its likelihood is obtained by marginalization:
\begin{align}
  p_\theta^{\mathrm{raw}}(w_{1:n-1},z_{1:n})
  =\int_0^1 f_\theta(\xi;w,z)\dd\xi.
  \label{eq:appendix-raw-marginal}
\end{align}
This integral must cover all state-implied reservoir fractions.  The
auxiliary density $q_\epsilon$ is supported only on
$I_\epsilon=[\epsilon_{\min},\epsilon_{\max}]$, so define the restricted
mass
\begin{equation}
  Z_{I_\epsilon}(w,z)
  :=\int_{I_\epsilon}f_\theta(\xi;w,z)\dd\xi.
\end{equation}
Insert $q_\epsilon(\xi)/q_\epsilon(\xi)$ only on this interval.  Then
\begin{align}
  p_\theta^{\mathrm{raw}}(w_{1:n-1},z_{1:n})
  \geq Z_{I_\epsilon}(w,z)
  =\E_{\xi\sim q_\epsilon}\left[
    \frac{f_\theta(\xi;w,z)}{q_\epsilon(\xi)}
  \right].
  \label{eq:appendix-importance-identity}
\end{align}
The first relation is an inequality unless the endpoint model has no
conditional mass outside $I_\epsilon$.  Monotonicity of the logarithm and
then Jensen's inequality yield
\begin{align}
  \log p_\theta^{\mathrm{raw}}(w_{1:n-1},z_{1:n})
  \geq\E_{\xi\sim q_\epsilon}\Big[
    &\log p_{\theta,n}(X(\xi))
    +n\log M
    +(n-1)\log(1-\xi)
    -\log q_\epsilon(\xi)
  \Big],
  \label{eq:appendix-dequantization-bound}
\end{align}
which is \eqref{eq:dequantization-bound}.

Assume the displayed raw likelihood and $Z_{I_\epsilon}(w,z)$ are positive.
The gap then has the usual variational interpretation.  Define
\begin{equation}
  p_\theta(\xi\mid w,z)
  =\frac{f_\theta(\xi;w,z)}
  {p_\theta^{\mathrm{raw}}(w_{1:n-1},z_{1:n})}.
  \label{eq:appendix-dequantization-posterior}
\end{equation}
Then the difference between the left side of
\eqref{eq:appendix-dequantization-bound} and its right side is
\begin{equation}
  \KL\!\left(q_\epsilon(\xi)\,\middle\|\,
  p_\theta(\xi\mid w,z)\right)\geq0.
  \label{eq:appendix-dequantization-gap}
\end{equation}
Equivalently, the same gap can be separated as
\begin{equation}
  \log\frac{p_\theta^{\mathrm{raw}}(w,z)}{Z_{I_\epsilon}(w,z)}
  +\KL\!\left(
    q_\epsilon(\xi)\,\middle\|\,
    \frac{f_\theta(\xi;w,z)}{Z_{I_\epsilon}(w,z)}
  \right).
  \label{eq:appendix-dequantization-gap-split}
\end{equation}
The bound is exact only when the full conditional posterior is supported on
$I_\epsilon$ and equals $q_\epsilon$ almost everywhere.

\section{Full proof of the split counting factor}
\label{app:counting}

We give a more explicit version of the calculation behind
Proposition~\ref{prop:reverse-split}.  Hold the $n-1$ unaffected components fixed and call their coordinate element $\dd O$.  Selecting one parent from an unordered $n$-component state changes the measure factor as
\begin{equation}
  \frac{1}{n!}\sum_{i=1}^{n}\dd O\dd m\dd z
  =\frac{1}{(n-1)!}\dd O\dd m\dd z.
  \label{eq:parent-counting}
\end{equation}
Adding the canonical split-mark element gives
\begin{equation}
  \frac{1}{(n-1)!}\dd O\dd m\dd z\dd\rho\dd u.
  \label{eq:forward-split-volume}
\end{equation}

In the $(n+1)$-component child state, select one physical unordered pair.
The labeled-coordinate integral contains both label assignments for its two
members, even though the sum below is written with $a<b$.  Counting the
candidate pairs gives
\begin{equation}
  \frac{1}{(n+1)!}\sum_{a<b}\dd O\dd m_a\dd m_b\dd z_a\dd z_b
  =\frac{1}{2(n-1)!}
  \dd O\dd m_a\dd m_b\dd z_a\dd z_b.
  \label{eq:pair-counting}
\end{equation}
Lemma~\ref{lem:split-jacobian} gives
\begin{equation}
  \dd m_a\dd m_b\dd z_a\dd z_b
  =m\dd m\dd\rho\dd z\dd u.
\end{equation}
The full labeled child integral in \eqref{eq:pair-counting} is twice the
integral over the canonical orientation: one labeled assignment has
$z_b-z_a\in\mathcal U_+$ and the other has its negative.  Rewriting the
displayed measure on the one-orientation canonical chart therefore cancels
the factor $1/2$.  On that chart, Lemma~\ref{lem:split-jacobian} gives
\begin{equation}
  \dd m_a\dd m_b\dd z_a\dd z_b
  =m\dd m\dd\rho\dd z\dd u,
\end{equation}
so \eqref{eq:forward-split-volume} equals $1/m$ times the canonically
oriented pair measure.  This produces \eqref{eq:exact-merge-rate}.

If both $u$ and $-u$ are retained as distinct split marks, the same physical
split has two coordinate representations and the corresponding density
convention must compensate for that duplication.  The final physical rate
is unchanged.  This is why the paper fixes a canonical half-space at the
outset.

\section{Forward path density and the fixed ELBO terms}
\label{app:forward-density}

Let split events occur at times $t_j^S$ with marks
$(i_j,\rho_j,u_j)$, and let deletion events occur at $t_k^D$ with
deleted component $i_k$.  Relative to the forward event times and split
marks, the forward log density is
\begin{align}
  \log q_X^F(H)
  ={}&
  -\int_0^\tau
    \ind\{N(X_t)<N_{\max}\}N(X_t)\beta(t)\dd t
  +\sum_j
  \left[
    \log\beta(t_j^S)
    +\log\kappa_{t_j^S}(\rho_j,u_j\mid X_{t_j^-},i_j)
  \right]
  \nonumber\\
  &-\int_\tau^T N(X_t)\delta(t)\dd t
  +\sum_k\log\delta(t_k^D).
  \label{eq:forward-path-density}
\end{align}
The needed change of variables is the joint map $\Psi:(X,H)\mapsto R$ from
\eqref{eq:joint-reversal-map}.  We now give its continuous determinant.  On
a forward flow segment with $n\geq1$ live components, let
$c=\sum_i m_i$ at the start and let $G$ be the integrated evaporation rate.
The segment maps the total condensed mass to
\begin{equation}
  f_G(c)=\frac{Mce^{-G}}{M-c+ce^{-G}},
  \qquad
  \alpha_G(c)=\frac{f_G(c)}{c}.
  \label{eq:flow-radial-map}
\end{equation}
Since $m_i'=\alpha_G(c)m_i$, the determinant on the $n$ mass coordinates is
\begin{equation}
  J_{\mathrm{flow}}(c,G,n)
  =\alpha_G(c)^{n-1}f_G'(c),
  \qquad
  f_G'(c)=
  \frac{M^2e^{-G}}{(M-c+ce^{-G})^2}.
  \label{eq:flow-jacobian}
\end{equation}
To see this directly, $n-1$ directions preserve $c$ and are multiplied by
$\alpha_G$; the remaining radial direction changes $c$ by $f_G'(c)$.

Fix a discrete lineage pattern.  Let $\mathcal I_F$ be its nonempty flow
segments, let $\mathcal E_S$ index its split events, and let
$m_j^{\mathrm{parent}}$ be the parent mass immediately before split $j$.
Composing the flow maps, split maps, deletions, and event-time reversal gives
\begin{equation}
  J_\Psi(X,H)
  =\prod_{\ell\in\mathcal I_F}
    \left[
      \alpha_{G_\ell}(c_\ell)^{n_\ell-1}
      f_{G_\ell}'(c_\ell)
    \right]
    \prod_{j\in\mathcal E_S}m_j^{\mathrm{parent}}.
  \label{eq:joint-history-jacobian}
\end{equation}
An empty flow segment contributes one.  Lemma~\ref{lem:split-jacobian}
gives each parent-mass factor.  Deletion merely moves one component's
coordinates into its future birth mark, so its determinant is one.
Reordering events and replacing $t$ by $T-t$ also have absolute determinant
one.  The derivative is block triangular when event times are retained as
coordinates, so their effect on later flow states does not alter the product
of diagonal blocks.  Canonical orientation removes the duplicate discrete
representation.
On each canonical lineage chart, this determinant is the coordinate form of
the Radon--Nikodym factor used in
Lemma~\ref{lem:joint-history-change}.

Equations \eqref{eq:forward-path-density} and
\eqref{eq:joint-history-jacobian}, together with any reservoir
dequantization density, supply all fixed terms in
Theorem~\ref{thm:path-elbo}.  They are not needed for parameter gradients but
must be included in a numerical ELBO.

Although $\delta(t)$ diverges, the expression is finite for every realized history because all event times are strictly below $T$, and after the last deletion $N(X_t)=0$.  Equivalently, one may evaluate the deletion contribution directly using the uniform deletion-time density $1/(T-\tau)$ for each component alive at $\tau$.

\section{Derivation of the point-process loss}
\label{app:loss-derivation}

Conditional on $L$, the ordered Stage-N birth times have joint density
$L!/S_N^L$.  Birth marks are sampled sequentially from
$A_\theta^{(m)}$.  In
Stage C, the usual marked point-process survival calculation gives the
reverse-record log density
\begin{align}
  \log\mathcal L_\theta(R)
  ={}&\log\Pi_\theta(L\mid M,y)
  +\log(L!)-L\log S_N
  +\sum_{e\in\mathcal E_B}
    \log A_\theta^{(m)}(m_e,z_e\mid X_{s_e^-},s_e,L)
  \nonumber\\
  &-\int_{S_N}^{T}K_\theta^{\mathrm{tot}}(X_s,s)\dd s
  +\sum_{e\in\mathcal E_C}
    \log K_\theta(a_e,b_e\mid X_{s_e^-},s_e).
  \label{eq:reverse-record-density}
\end{align}
The negative of the $\theta$-dependent terms is exactly
\eqref{eq:path-loss}.  The fixed birth-time density explains why no learned
Stage-N survival integral appears there.

\section{A finite-state validation model}
\label{app:finite-validation}

The exact experiment in Section~\ref{sec:experiments} can use the following finite analogue.

Let total mass be an integer $M_{\mathrm{int}}$, features lie on a lattice
$\{1,\ldots,K_{\mathrm{lat}}\}$, and component masses are positive
integers.  Restrict cardinality to $N_{\max}$ and use only discrete split
marks that keep masses and features on these lattices.  Enumerate every
unordered multiset
\begin{equation}
  X=\left(v;\{(m_i,z_i)\}_{i=1}^{n}\right)
\end{equation}
with integer entries satisfying the mass constraint.  Set $\gamma\equiv0$
in this validation toy, because the continuous evaporation ODE does not
preserve integer masses.  Define the \emph{forward} generator $G_t$ using
only the discrete split and deletion maps.  Births and merges will appear in
its time reversal rather than being added to the forward generator.

Use the row-vector convention for the marginal probability vector $P_t$:
\begin{equation}
  \frac{\dd P_t}{\dd t}=P_tG_t.
  \label{eq:finite-forward-equation}
\end{equation}
This finite ODE can be integrated directly for a time-dependent schedule.
On an interval where $G_t$ is constant, its exact update is
$P_{t+\Delta}=P_t\exp(\Delta G_t)$.  The singular deletion schedule is not
piecewise constant, so it should be handled by the ODE or its known
integrated hazard, not by one constant-matrix exponential.  Compute reverse
rates for $t<T$ and take the all-reservoir terminal law as the limit
$t\uparrow T$.
The exact reverse off-diagonal generator is
\begin{equation}
  \overline G_{T-t}(Y,X)
  =G_t(X,Y)\frac{P_t(X)}{P_t(Y)},
  \qquad X\neq Y, P_t(Y)>0,
  \label{eq:finite-reverse-generator}
\end{equation}
and diagonal entries make rows sum to zero.  This benchmark can verify flux
reversal, pair counting, the training objective, endpoint distributions, and
conservation without numerical density estimation.  Because it uses counting
measure, it cannot validate the continuous $1/m$ Jacobian.  That factor is
checked separately by sampling smooth test functions on both sides of the
change-of-variables identity in Lemma~\ref{lem:split-jacobian} and comparing
Monte Carlo integrals.

\section{Implementation checklist}
\label{app:implementation-checklist}

Before running large experiments, an implementation should pass the following tests.

\begin{enumerate}
  \item \textbf{Map inversion:} randomly split a component and merge its children; recover the original state to numerical tolerance.
  \item \textbf{Mass:} after every flow and event, verify
  $|v+\sum_i m_i-M|/M<10^{-10}$ in double precision.
  \item \textbf{Permutation:} permute input components and verify invariant total rates and correspondingly permuted pair probabilities.
  \item \textbf{Terminal state:} simulate many forward paths and verify that all are empty at $T$.
  \item \textbf{Point-process likelihood:} on a constant-rate toy process, compare the implemented path likelihood with the analytic exponential waiting-time likelihood.
  \item \textbf{Integral estimator:} compare Monte Carlo survival integrals with dense numerical quadrature on stored trajectories.
  \item \textbf{Coordinate Jacobians:} compare automatic or finite-difference
  determinants with \eqref{eq:dequantization-jacobian},
  \eqref{eq:split-jacobian}, and \eqref{eq:flow-jacobian}; separately verify
  the continuous $1/m$ change of variables by Monte Carlo integration.
  \item \textbf{Birth bridge:} conditional on each tested $L$, verify that
  exactly $L$ births occur and that their ordered times match the order
  statistics of $L$ independent uniform variables.
  \item \textbf{Finite reverse:} compare learned and analytical reverse generators in Appendix~\ref{app:finite-validation}.
  \item \textbf{Thinning:} verify every proposed envelope exceeds the evaluated rate.  An invalid envelope biases samples.
  \item \textbf{Genealogy bookkeeping:} reverse every simulated forward history and check that it returns exactly to $X_0$ before any neural model is introduced.
\end{enumerate}

\begin{thebibliography}{99}

\bibitem{aldous1999deterministic}
D.~J. Aldous.
\newblock Deterministic and stochastic models for coalescence (aggregation and coagulation): A review of the mean-field theory for probabilists.
\newblock \emph{Bernoulli}, 5(1):3--48, 1999.

\bibitem{anderson1982reverse}
B.~D.~O. Anderson.
\newblock Reverse-time diffusion equation models.
\newblock \emph{Stochastic Processes and their Applications}, 12(3):313--326, 1982.

\bibitem{andreassen2019junipr}
A.~Andreassen, I.~Feige, C.~Frye, and M.~D. Schwartz.
\newblock JUNIPR: A framework for unsupervised machine learning in particle physics.
\newblock \emph{European Physical Journal C}, 79:102, 2019.

\bibitem{bansal2023cold}
A.~Bansal, E.~Borgnia, H.-M.~Chu, J.~Li, H.~Kazemi, F.~Huang, M.~Goldblum, J.~Geiping, and T.~Goldstein.
\newblock Cold Diffusion: Inverting arbitrary image transforms without noise.
\newblock \emph{Advances in Neural Information Processing Systems}, 36, 2023.

\bibitem{bergmeister2024efficient}
A.~Bergmeister, K.~Martinkus, N.~Perraudin, and R.~Wattenhofer.
\newblock Efficient and scalable graph generation through iterative local expansion.
\newblock In \emph{International Conference on Learning Representations}, 2024.

\bibitem{bertoin2006random}
J.~Bertoin.
\newblock \emph{Random Fragmentation and Coagulation Processes}.
\newblock Cambridge University Press, 2006.

\bibitem{boget2026hierarchical}
Y.~Boget, P.~Strasser, and A.~Kalousis.
\newblock Hierarchical discrete flow matching for graph generation.
\newblock arXiv:2604.00236, 2026.

\bibitem{campbell2023transdimensional}
A.~Campbell, W.~Harvey, C.~Weilbach, V.~De~Bortoli, T.~Rainforth, and A.~Doucet.
\newblock Trans-dimensional generative modeling via jump diffusion models.
\newblock \emph{Advances in Neural Information Processing Systems}, 36, 2023.

\bibitem{chen2023learning}
T.~Chen and M.~Zhou.
\newblock Learning to Jump: Thinning and thickening latent counts for generative modeling.
\newblock In \emph{Proceedings of the 40th International Conference on Machine Learning}, 2023.

\bibitem{conforti2022time}
G.~Conforti and C.~L\'eonard.
\newblock Time reversal of Markov processes with jumps under a finite entropy condition.
\newblock \emph{Stochastic Processes and their Applications}, 144:85--124, 2022.

\bibitem{green1995reversible}
P.~J. Green.
\newblock Reversible jump Markov chain Monte Carlo computation and Bayesian model determination.
\newblock \emph{Biometrika}, 82(4):711--732, 1995.

\bibitem{ho2020denoising}
J.~Ho, A.~Jain, and P.~Abbeel.
\newblock Denoising diffusion probabilistic models.
\newblock \emph{Advances in Neural Information Processing Systems}, 33, 2020.

\bibitem{lee2019settransformer}
J.~Lee, Y.~Lee, J.~Kim, A.~Kosiorek, S.~Choi, and Y.~W. Teh.
\newblock Set Transformer: A framework for attention-based permutation-invariant neural networks.
\newblock In \emph{Proceedings of the 36th International Conference on Machine Learning}, 2019.

\bibitem{norris1999smoluchowski}
J.~R. Norris.
\newblock Smoluchowski's coagulation equation: Uniqueness, nonuniqueness and a hydrodynamic limit for the stochastic coalescent.
\newblock \emph{Annals of Applied Probability}, 9(1):78--109, 1999.

\bibitem{pan2026existence}
X.~Pan, C.~R. Shelton, R.~Mahishi, and C.~Hong.
\newblock Existence-Field Diffusion Model for spatial point processes with variable cardinality.
\newblock arXiv:2607.26428, 2026.

\bibitem{sohldickstein2015deep}
J.~Sohl-Dickstein, E.~Weiss, N.~Maheswaranathan, and S.~Ganguli.
\newblock Deep unsupervised learning using nonequilibrium thermodynamics.
\newblock In \emph{Proceedings of the 32nd International Conference on Machine Learning}, 2015.

\bibitem{song2021score}
Y.~Song, J.~Sohl-Dickstein, D.~P. Kingma, A.~Kumar, S.~Ermon, and B.~Poole.
\newblock Score-based generative modeling through stochastic differential equations.
\newblock In \emph{International Conference on Learning Representations}, 2021.

\bibitem{zaheer2017deepsets}
M.~Zaheer, S.~Kottur, S.~Ravanbakhsh, B.~Poczos, R.~Salakhutdinov, and A.~Smola.
\newblock Deep Sets.
\newblock \emph{Advances in Neural Information Processing Systems}, 30, 2017.

\end{thebibliography}

\end{document}
